% Options for packages loaded elsewhere
%DIF LATEXDIFF DIFFERENCE FILE
%DIF DEL peer_review/submissions/2025-06-07-PNAS/Manuscript/main_20250607_5c14834.tex   Wed Dec  3 16:40:48 2025
%DIF ADD docs/report.tex                                                                Fri Dec 12 13:39:53 2025
\PassOptionsToPackage{unicode}{hyperref}
\PassOptionsToPackage{hyphens}{url}
\PassOptionsToPackage{dvipsnames,svgnames,x11names}{xcolor}
%
\documentclass[
%DIF 7c7
%DIF <   12pt,
%DIF -------
  11pt, %DIF > 
%DIF -------
]{article}
\usepackage{amsmath,amssymb}
\usepackage{setspace}
\usepackage{iftex}
\ifPDFTeX
  \usepackage[T1]{fontenc}
  \usepackage[utf8]{inputenc}
  \usepackage{textcomp} % provide euro and other symbols
\else % if luatex or xetex
  \usepackage{unicode-math} % this also loads fontspec
  \defaultfontfeatures{Scale=MatchLowercase}
  \defaultfontfeatures[\rmfamily]{Ligatures=TeX,Scale=1}
\fi
\usepackage{lmodern}
\ifPDFTeX\else
  % xetex/luatex font selection
    \setmainfont[]{Arial}
\fi
% Use upquote if available, for straight quotes in verbatim environments
\IfFileExists{upquote.sty}{\usepackage{upquote}}{}
\IfFileExists{microtype.sty}{% use microtype if available
  \usepackage[]{microtype}
  \UseMicrotypeSet[protrusion]{basicmath} % disable protrusion for tt fonts
}{}
\makeatletter
\@ifundefined{KOMAClassName}{% if non-KOMA class
  \IfFileExists{parskip.sty}{%
    \usepackage{parskip}
  }{% else
    \setlength{\parindent}{0pt}
    \setlength{\parskip}{6pt plus 2pt minus 1pt}}
}{% if KOMA class
  \KOMAoptions{parskip=half}}
\makeatother
\usepackage{xcolor}
\usepackage[margin=1in]{geometry}
\usepackage{longtable,booktabs,array}
\usepackage{calc} % for calculating minipage widths
% Correct order of tables after \paragraph or \subparagraph
\usepackage{etoolbox}
\makeatletter
\patchcmd\longtable{\par}{\if@noskipsec\mbox{}\fi\par}{}{}
\makeatother
% Allow footnotes in longtable head/foot
\IfFileExists{footnotehyper.sty}{\usepackage{footnotehyper}}{\usepackage{footnote}}
\makesavenoteenv{longtable}
\usepackage{graphicx}
\makeatletter
\newsavebox\pandoc@box
\newcommand*\pandocbounded[1]{% scales image to fit in text height/width
  \sbox\pandoc@box{#1}%
  \Gscale@div\@tempa{\textheight}{\dimexpr\ht\pandoc@box+\dp\pandoc@box\relax}%
  \Gscale@div\@tempb{\linewidth}{\wd\pandoc@box}%
  \ifdim\@tempb\p@<\@tempa\p@\let\@tempa\@tempb\fi% select the smaller of both
  \ifdim\@tempa\p@<\p@\scalebox{\@tempa}{\usebox\pandoc@box}%
  \else\usebox{\pandoc@box}%
  \fi%
}
% Set default figure placement to htbp
\def\fps@figure{htbp}
\makeatother
\setlength{\emergencystretch}{3em} % prevent overfull lines
\providecommand{\tightlist}{%
  \setlength{\itemsep}{0pt}\setlength{\parskip}{0pt}}
\setcounter{secnumdepth}{5}
% definitions for citeproc citations
\NewDocumentCommand\citeproctext{}{}
\NewDocumentCommand\citeproc{mm}{%
  \begingroup\def\citeproctext{#2}\cite{#1}\endgroup}
\makeatletter
 % allow citations to break across lines
 \let\@cite@ofmt\@firstofone
 % avoid brackets around text for \cite:
 \def\@biblabel#1{}
 \def\@cite#1#2{{#1\if@tempswa , #2\fi}}
\makeatother
\newlength{\cslhangindent}
\setlength{\cslhangindent}{1.5em}
\newlength{\csllabelwidth}
\setlength{\csllabelwidth}{3em}
\newenvironment{CSLReferences}[2] % #1 hanging-indent, #2 entry-spacing
 {\begin{list}{}{%
  \setlength{\itemindent}{0pt}
  \setlength{\leftmargin}{0pt}
  \setlength{\parsep}{0pt}
  % turn on hanging indent if param 1 is 1
  \ifodd #1
   \setlength{\leftmargin}{\cslhangindent}
   \setlength{\itemindent}{-1\cslhangindent}
  \fi
  % set entry spacing
  \setlength{\itemsep}{#2\baselineskip}}}
 {\end{list}}
\usepackage{calc}
\newcommand{\CSLBlock}[1]{\hfill\break\parbox[t]{\linewidth}{\strut\ignorespaces#1\strut}}
\newcommand{\CSLLeftMargin}[1]{\parbox[t]{\csllabelwidth}{\strut#1\strut}}
\newcommand{\CSLRightInline}[1]{\parbox[t]{\linewidth - \csllabelwidth}{\strut#1\strut}}
\newcommand{\CSLIndent}[1]{\hspace{\cslhangindent}#1}
\ifLuaTeX
\usepackage[bidi=basic]{babel}
\else
\usepackage[bidi=default]{babel}
\fi
\babelprovide[main,import]{british}
\ifPDFTeX
\else
\babelfont{rm}[]{Arial}
\fi
% get rid of language-specific shorthands (see #6817):
\let\LanguageShortHands\languageshorthands
\def\languageshorthands#1{}
\ifLuaTeX
  \usepackage[english]{selnolig} % disable illegal ligatures
\fi
%DIF 122-123d122
%DIF < %\usepackage{times}
%DIF < \usepackage{babel}
%DIF -------
\usepackage{colortbl}
\usepackage[footnotesize,bf,singlelinecheck=false]{caption}
%DIF 126-134d124
%DIF < 
%DIF < \usepackage{hyperref}
%DIF < \hypersetup{
%DIF <   colorlinks=true,
%DIF <   linkcolor=blue,
%DIF <   citecolor=blue,
%DIF <   urlcolor=blue
%DIF < }
%DIF < % for drafting
%DIF -------
\usepackage[nomarkers,nolists]{endfloat} % nomarkers, nolists
%DIF 136-137c125-126
%DIF < \let\oldsection\section
%DIF < \renewcommand{\section}{\clearpage\oldsection}
%DIF -------
%\let\oldsection\section %DIF > 
%\renewcommand{\section}{\clearpage\oldsection} %DIF > 
%DIF -------

\usepackage{booktabs}
\usepackage{longtable}
\usepackage{array}
\usepackage{multirow}
\usepackage{wrapfig}
\usepackage{float}
\usepackage{colortbl}
\usepackage{pdflscape}
\usepackage{tabu}
\usepackage{threeparttable}
\usepackage{threeparttablex}
\usepackage[normalem]{ulem}
\usepackage{makecell}
\usepackage{xcolor}
\usepackage{bookmark}
\IfFileExists{xurl.sty}{\usepackage{xurl}}{} % add URL line breaks if available
\urlstyle{same}
\hypersetup{
  pdftitle={Gender in Teaching: Insights from Five Million Syllabi on Collaboration, Interdisciplinarity, and Reading Selections},
  pdfauthor={Andrea Blasco; Elisa Orlandi},
  pdflang={en-UK},
  colorlinks=true,
  linkcolor={Maroon},
  filecolor={Maroon},
  citecolor={Blue},
  urlcolor={Blue},
  pdfcreator={LaTeX via pandoc}}

\title{Gender in Teaching:
Insights from Five Million Syllabi on Collaboration, Interdisciplinarity, and Reading Selections}
\author{Andrea Blasco\footnote{\DIFaddbegin \DIFadd{Corresponding author; }\DIFaddend European Commission, Joint Research Centre, Brussels, Belgium; \DIFdelbegin \DIFdel{(}\DIFdelend \href{mailto:andrea.blasco@ec.europa.eu}{\nolinkurl{andrea.blasco@ec.europa.eu}}\DIFdelbegin \DIFdel{)}\DIFdelend } \and Elisa Orlandi\DIFdelbegin \footnote{\DIFdel{University of Bologna, Bologna, Italy, and European Commission, Joint Research Centre, Brussels, Belgium; (}\href{mailto:elisa.orlandi15@unibo.it}{%DIFDELCMD < \nolinkurl{elisa.orlandi15@unibo.it}%%%
}%DIFAUXCMD
\DIFdel{)}}%DIFAUXCMD
\DIFdelend \DIFaddbegin \footnote{\DIFadd{University of Salerno, Salerno, Italy, and European Commission, Joint Research Centre, Brussels, Belgium; }\href{mailto:eoralndi@unisa.it}{\nolinkurl{eoralndi@unisa.it}}}\DIFaddend }
\date{}
%DIF PREAMBLE EXTENSION ADDED BY LATEXDIFF
%DIF UNDERLINE PREAMBLE %DIF PREAMBLE
\RequirePackage[normalem]{ulem} %DIF PREAMBLE
\RequirePackage{color}\definecolor{RED}{rgb}{1,0,0}\definecolor{BLUE}{rgb}{0,0,1} %DIF PREAMBLE
\providecommand{\DIFadd}[1]{{\protect\color{blue}\uwave{#1}}} %DIF PREAMBLE
\providecommand{\DIFdel}[1]{{\protect\color{red}\sout{#1}}}                      %DIF PREAMBLE
%DIF SAFE PREAMBLE %DIF PREAMBLE
\providecommand{\DIFaddbegin}{} %DIF PREAMBLE
\providecommand{\DIFaddend}{} %DIF PREAMBLE
\providecommand{\DIFdelbegin}{} %DIF PREAMBLE
\providecommand{\DIFdelend}{} %DIF PREAMBLE
\providecommand{\DIFmodbegin}{} %DIF PREAMBLE
\providecommand{\DIFmodend}{} %DIF PREAMBLE
%DIF FLOATSAFE PREAMBLE %DIF PREAMBLE
\providecommand{\DIFaddFL}[1]{\DIFadd{#1}} %DIF PREAMBLE
\providecommand{\DIFdelFL}[1]{\DIFdel{#1}} %DIF PREAMBLE
\providecommand{\DIFaddbeginFL}{} %DIF PREAMBLE
\providecommand{\DIFaddendFL}{} %DIF PREAMBLE
\providecommand{\DIFdelbeginFL}{} %DIF PREAMBLE
\providecommand{\DIFdelendFL}{} %DIF PREAMBLE
\newcommand{\DIFscaledelfig}{0.5}
%DIF HIGHLIGHTGRAPHICS PREAMBLE %DIF PREAMBLE
\RequirePackage{settobox} %DIF PREAMBLE
\RequirePackage{letltxmacro} %DIF PREAMBLE
\newsavebox{\DIFdelgraphicsbox} %DIF PREAMBLE
\newlength{\DIFdelgraphicswidth} %DIF PREAMBLE
\newlength{\DIFdelgraphicsheight} %DIF PREAMBLE
% store original definition of \includegraphics %DIF PREAMBLE
\LetLtxMacro{\DIFOincludegraphics}{\includegraphics} %DIF PREAMBLE
\newcommand{\DIFaddincludegraphics}[2][]{{\color{blue}\fbox{\DIFOincludegraphics[#1]{#2}}}} %DIF PREAMBLE
\newcommand{\DIFdelincludegraphics}[2][]{% %DIF PREAMBLE
\sbox{\DIFdelgraphicsbox}{\DIFOincludegraphics[#1]{#2}}% %DIF PREAMBLE
\settoboxwidth{\DIFdelgraphicswidth}{\DIFdelgraphicsbox} %DIF PREAMBLE
\settoboxtotalheight{\DIFdelgraphicsheight}{\DIFdelgraphicsbox} %DIF PREAMBLE
\scalebox{\DIFscaledelfig}{% %DIF PREAMBLE
\parbox[b]{\DIFdelgraphicswidth}{\usebox{\DIFdelgraphicsbox}\\[-\baselineskip] \rule{\DIFdelgraphicswidth}{0em}}\llap{\resizebox{\DIFdelgraphicswidth}{\DIFdelgraphicsheight}{% %DIF PREAMBLE
\setlength{\unitlength}{\DIFdelgraphicswidth}% %DIF PREAMBLE
\begin{picture}(1,1)% %DIF PREAMBLE
\thicklines\linethickness{2pt} %DIF PREAMBLE
{\color[rgb]{1,0,0}\put(0,0){\framebox(1,1){}}}% %DIF PREAMBLE
{\color[rgb]{1,0,0}\put(0,0){\line( 1,1){1}}}% %DIF PREAMBLE
{\color[rgb]{1,0,0}\put(0,1){\line(1,-1){1}}}% %DIF PREAMBLE
\end{picture}% %DIF PREAMBLE
}\hspace*{3pt}}} %DIF PREAMBLE
} %DIF PREAMBLE
\LetLtxMacro{\DIFOaddbegin}{\DIFaddbegin} %DIF PREAMBLE
\LetLtxMacro{\DIFOaddend}{\DIFaddend} %DIF PREAMBLE
\LetLtxMacro{\DIFOdelbegin}{\DIFdelbegin} %DIF PREAMBLE
\LetLtxMacro{\DIFOdelend}{\DIFdelend} %DIF PREAMBLE
\DeclareRobustCommand{\DIFaddbegin}{\DIFOaddbegin \let\includegraphics\DIFaddincludegraphics} %DIF PREAMBLE
\DeclareRobustCommand{\DIFaddend}{\DIFOaddend \let\includegraphics\DIFOincludegraphics} %DIF PREAMBLE
\DeclareRobustCommand{\DIFdelbegin}{\DIFOdelbegin \let\includegraphics\DIFdelincludegraphics} %DIF PREAMBLE
\DeclareRobustCommand{\DIFdelend}{\DIFOaddend \let\includegraphics\DIFOincludegraphics} %DIF PREAMBLE
\LetLtxMacro{\DIFOaddbeginFL}{\DIFaddbeginFL} %DIF PREAMBLE
\LetLtxMacro{\DIFOaddendFL}{\DIFaddendFL} %DIF PREAMBLE
\LetLtxMacro{\DIFOdelbeginFL}{\DIFdelbeginFL} %DIF PREAMBLE
\LetLtxMacro{\DIFOdelendFL}{\DIFdelendFL} %DIF PREAMBLE
\DeclareRobustCommand{\DIFaddbeginFL}{\DIFOaddbeginFL \let\includegraphics\DIFaddincludegraphics} %DIF PREAMBLE
\DeclareRobustCommand{\DIFaddendFL}{\DIFOaddendFL \let\includegraphics\DIFOincludegraphics} %DIF PREAMBLE
\DeclareRobustCommand{\DIFdelbeginFL}{\DIFOdelbeginFL \let\includegraphics\DIFdelincludegraphics} %DIF PREAMBLE
\DeclareRobustCommand{\DIFdelendFL}{\DIFOaddendFL \let\includegraphics\DIFOincludegraphics} %DIF PREAMBLE
%DIF COLORLISTINGS PREAMBLE %DIF PREAMBLE
\RequirePackage{listings} %DIF PREAMBLE
\RequirePackage{color} %DIF PREAMBLE
\lstdefinelanguage{DIFcode}{ %DIF PREAMBLE
%DIF DIFCODE_UNDERLINE %DIF PREAMBLE
  moredelim=[il][\color{red}\sout]{\%DIF\ <\ }, %DIF PREAMBLE
  moredelim=[il][\color{blue}\uwave]{\%DIF\ >\ } %DIF PREAMBLE
} %DIF PREAMBLE
\lstdefinestyle{DIFverbatimstyle}{ %DIF PREAMBLE
	language=DIFcode, %DIF PREAMBLE
	basicstyle=\ttfamily, %DIF PREAMBLE
	columns=fullflexible, %DIF PREAMBLE
	keepspaces=true %DIF PREAMBLE
} %DIF PREAMBLE
\lstnewenvironment{DIFverbatim}{\lstset{style=DIFverbatimstyle}}{} %DIF PREAMBLE
\lstnewenvironment{DIFverbatim*}{\lstset{style=DIFverbatimstyle,showspaces=true}}{} %DIF PREAMBLE
%DIF END PREAMBLE EXTENSION ADDED BY LATEXDIFF

\begin{document}
\maketitle
\begin{abstract}
This study examines the formation of academic teaching teams and how their gender composition is associated with course design and gender disparity in academic recognition. Analysing over five million syllabi from higher education in fifteen countries, we document a persistent underrepresentation of mixed-gender co-teaching teams, which occur 50\% less frequently than expected under gender-neutral team formation. We link team gender configurations to critical course characteristics: interdisciplinarity, novelty of reading materials, and gender balance in assigned citations. Courses taught by mixed-gender teams are significantly more interdisciplinary than those taught by same-gender teams. Male-only teams cite substantially fewer female authors than female-only teams, and mixed-gender teams do not eliminate this gap. Moreover, courses led by solo female instructors or teams including at least one woman tend to incorporate more \DIFdelbegin \DIFdel{novel reading materials}\DIFdelend \DIFaddbegin \DIFadd{recent reading materials, while mixed-gender teams tend to cite more conventional readings and fewer readings at the research frontier}\DIFaddend . These findings reveal persistent gendered patterns in teaching collaboration and content and point to mixed-gender co-teaching as a potential lever to promote interdisciplinarity and gender equity in academic representation.
\end{abstract}

\setstretch{1}
JEL Classification: I21, I23, J16

Keywords: education \textbar{} gender bias \textbar{} teams \textbar{} interdisciplinarity \textbar{} novelty \textbar{} syllabi

\newpage
\DIFdelbegin %DIFDELCMD < \doublespace
%DIFDELCMD < %%%
\DIFdelend 

\DIFaddbegin \section*{\DIFadd{Significance Statement}}\label{significance-statement}
\addcontentsline{toc}{section}{\DIFadd{Significance Statement}}

\DIFadd{Despite growing gender diversity in academia, we reveal a persistent underrepresentation of mixed-gender teaching teams across universities. By analysing 5M syllabi (1999-2019), we find that instructors tend to collaborate with same-gender colleagues at a rate more than twice the expected rate. This is not fully explained by institutional or disciplinary constraints, suggesting a gender-based bias. Moreover, mixed-gender teams tend to design more interdisciplinary courses and assign more gender-diverse readings, whereas women tend to assign more recent readings. These findings extend the literature on gender bias in academia into the domain of teaching, revealing how collaboration dynamics influence the pedagogical content and the student experience. Understanding these patterns can inform policies aimed at improving equity and pedagogical quality in higher education.
}

\DIFaddend \section{Introduction}\label{introduction}

Universities \DIFdelbegin \DIFdel{, vocational schools, }\DIFdelend and other higher education institutions \DIFdelbegin \DIFdel{play a crucial role in disseminating knowledgeand developing human capital . Instructors shape this role through their choices in teaching subjects, assigned readings, and co-teaching arrangements. These decisions have a significant impact on student }\DIFdelend \DIFaddbegin \DIFadd{are central to producing and disseminating knowledge, contributing to human capital development. Instructors' choices -- how to teach, which materials to assign, and whether to collaborate with colleagues -- shape students' educational experience, influencing their }\DIFaddend learning outcomes and \DIFdelbegin \DIFdel{students' career trajectories }\DIFdelend \DIFaddbegin \DIFadd{career trajectories substantially }\DIFaddend (\citeproc{ref-bakkenes2010teacher}{1}--\citeproc{ref-han2023interdisciplinary}{3}). \DIFdelbegin \DIFdel{However, they are not always optimal. Instructors may prefer working alone or }\DIFdelend \DIFaddbegin \DIFadd{These choices reflect both individual preferences and institutional constraints, which may not always align with students' learning needs or societal objectives. For instance, instructors may prefer to work alone or collaborate with colleagues }\DIFaddend with \DIFdelbegin \DIFdel{colleagues from }\DIFdelend similar backgrounds, which can limit course diversity \DIFdelbegin \DIFdel{and interdisciplinarity , }\DIFdelend \DIFaddbegin \DIFadd{or interdisciplinarity and }\DIFaddend reinforce gender stereotypes in specific subjects (\citeproc{ref-ellemers2018gender}{4}--\citeproc{ref-canaan2023impact}{6})\DIFdelbegin \DIFdel{, or even restrict instructors' exposure to different teaching methodologies, hindering their professional growth }\DIFdelend \DIFaddbegin \DIFadd{. Collaboration can also advance instructors' professional development by exposing them to different pedagogical approaches or teaching methods }\DIFaddend (\citeproc{ref-buckingham2021learning}{7}). \DIFdelbegin \DIFdel{Despite its significance, research on the underlying factors that drive }\DIFdelend \DIFaddbegin \DIFadd{Yet, despite their significance, the factors shaping }\DIFaddend these teaching choices \DIFdelbegin \DIFdel{remains limited}\DIFdelend \DIFaddbegin \DIFadd{remain poorly understood}\DIFaddend , primarily due to a lack of data.

This study addresses this gap by analysing a large dataset of \DIFdelbegin \DIFdel{English }\DIFdelend \DIFaddbegin \DIFadd{English-language }\DIFaddend course syllabi from 1990 to 2019 across fifteen countries \DIFdelbegin \DIFdel{, which has }\DIFdelend \DIFaddbegin \DIFadd{and from 2019 to 2023 across thirty-three European countries. These data, which have }\DIFaddend been previously used \DIFdelbegin \DIFdel{to study }\DIFdelend \DIFaddbegin \DIFadd{in research on }\DIFaddend school curricula (\citeproc{ref-han2023interdisciplinary}{3}), teaching metrics (\citeproc{ref-nowogrodzki2016mining}{8}), the impacts of AI (\citeproc{ref-gao2023quantifying}{9}), and the education-innovation gap (\citeproc{ref-biasi2022education}{10})\DIFdelbegin \DIFdel{. The dataset enables }\DIFdelend \DIFaddbegin \DIFadd{, enable }\DIFaddend a systematic, quantitative analysis of long-term trends. Each syllabus includes \DIFdelbegin \DIFdel{instructor }\DIFdelend \DIFaddbegin \DIFadd{the instructor's }\DIFaddend gender, institutional affiliation, course field, assigned readings, and a course description. This \DIFaddbegin \DIFadd{information }\DIFaddend allows us to examine gender-based team formation, revealing potential preferences and institutional constraints in teaching collaborations. Furthermore, we analyse how team composition, particularly gender and size, relates to teaching practices, such as interdisciplinarity and reading assignments, while controlling for institutional factors, field, academic year, and other relevant variables.

\DIFdelbegin \DIFdel{Specifically, this study addresses two key questions. First, how often do academics choose to teach alone versus co-teach with a colleague, and to what extent are co-teaching teams composed of mixed-gender pairs? Second, how do the size and gender composition }\DIFdelend \DIFaddbegin \DIFadd{We investigate three central questions: (i) how gender and institutional constraints are associated with the formation }\DIFaddend of teaching teams\DIFdelbegin \DIFdel{relate to key aspects of teaching: (1) the course's interdisciplinarity, (2) the novelty of assigned readings, and (3) the gender representation of the cited authors? Previous research has shown that these aspects }\DIFdelend \DIFaddbegin \DIFadd{; (ii) whether instructor gender is associated with the inclusion of female-authored readings; and (iii) whether mixed-gender teams are linked to greater interdisciplinarity in course content and with the adoption of novel or unconventional readings, including materials that are recent or at the frontier of research (``scientific novelty''). Previous research shows that these teaching choices }\DIFaddend shape students' learning outcomes and career trajectories (\citeproc{ref-bakkenes2010teacher}{1}--\citeproc{ref-han2023interdisciplinary}{3})\DIFdelbegin \DIFdel{while also influencing }\DIFdelend \DIFaddbegin \DIFadd{, and they can also influence }\DIFaddend instructors' professional development and promotion within their institutions (\citeproc{ref-buckingham2021learning}{7}).

\DIFaddbegin \DIFadd{We complement these analyses with an exploration of the formation of teaching teams in Europe pre- and post-COVID-19. The COVID-19 pandemic disrupted academic work in many countries, especially for women and caregivers (e.g., \mbox{%DIFAUXCMD
\citeproc{ref-myers2020unequal}{11}}\hskip0pt%DIFAUXCMD
), reshaped collaboration networks, and accelerated the shift to online and hybrid teaching. The shift to online and hybrid teaching may also have affected team formation across gender and disciplines. Using COVID-19 as an exogenous shock, we study changes in the prevalence and gender composition of teaching teams around 2020.
}

\DIFaddend Our results reveal a significant role for gender in co-teaching practices, with mixed-gender teams occurring consistently less frequently than same-gender teams across institutions and fields. Furthermore, mixed-gender teams arise much less often than expected under a null model that forms teams in a gender-neutral manner while \DIFdelbegin \DIFdel{keeping fixed }\DIFdelend \DIFaddbegin \DIFadd{holding }\DIFaddend field-related and institutional constraints \DIFaddbegin \DIFadd{fixed}\DIFaddend . These findings underscore a persistent and widespread underrepresentation of mixed-gender collaborations in teaching \DIFdelbegin \DIFdel{, that is not fully }\DIFdelend \DIFaddbegin \DIFadd{not }\DIFaddend explained by unobserved differences across fields or institutions.

Our analysis further reveals a significant and strong association between the gender composition of a teaching team and the course's interdisciplinarity, with mixed-gender teams being more interdisciplinary than all-male teams or courses taught by individual instructors. This evidence suggests that gender diversity can be a potential lever for increasing interdisciplinary teaching in tertiary education.
\DIFaddbegin 

\DIFaddend We also find a significant association between team configurations and the fraction of cited female authors. Specifically, classes taught by a female instructor alone tend to cite a higher fraction of female authors \DIFdelbegin \DIFdel{compared to courses }\DIFdelend \DIFaddbegin \DIFadd{than those }\DIFaddend taught by male instructors, with mixed-gender teams falling in between. This \DIFaddbegin \DIFadd{finding }\DIFaddend highlights teaching practices as an understudied potential driver of \DIFdelbegin \DIFdel{the underrepresentation of women}\DIFdelend \DIFaddbegin \DIFadd{women's underrepresentation }\DIFaddend in academia. Finally, we \DIFdelbegin \DIFdel{also }\DIFdelend observe a trend \DIFdelbegin \DIFdel{in novelty, with female instructors assigning }\DIFdelend \DIFaddbegin \DIFadd{toward novelty: female instructors assign }\DIFaddend more recent readings than their male counterparts, regardless of team size\DIFaddbegin \DIFadd{, but mixed-gender teams tend to cite more conventional readings and are less likely to cite cutting-edge research}\DIFaddend . This evidence contributes to the debate on the factors shaping the education-innovation gap (\citeproc{ref-biasi2022education}{10}).

\DIFdelbegin \DIFdel{While prior research has investigated academics decision-making, especially how they form research collaborations (\mbox{%DIFAUXCMD
\citeproc{ref-guimera2005team}{11}}\hskip0pt%DIFAUXCMD
}\DIFdelend \DIFaddbegin \section{\DIFadd{Literature}}\label{literature}

\DIFadd{Research shows that gender and institutional structures shape key aspects of academic work, from team formation to curriculum design. Four strands of this literature are particularly relevant to our study: (i) research on how gender influences collaboration patterns, (ii) the representation of women in course materials, (iii) the interdisciplinarity of mixed-gender teams, and (iv) the propensity of such teams to innovate. These insights provide the conceptual foundation for our study, which we review in detail below.
}

\subsection{\DIFadd{Gender and team formation in academia}}\label{gender-and-team-formation-in-academia}

\DIFadd{Teams play an increasingly important role in knowledge production (\mbox{%DIFAUXCMD
\citeproc{ref-wuchty2007increasing}{12}}\hskip0pt%DIFAUXCMD
, \mbox{%DIFAUXCMD
\citeproc{ref-milojevic2014principles}{13}}\hskip0pt%DIFAUXCMD
), and a large body of literature documents strong }\emph{\DIFadd{gender homophily}} \DIFadd{in scientific collaboration and professional networks (\mbox{%DIFAUXCMD
\citeproc{ref-mcpherson2001birds}{14}}\hskip0pt%DIFAUXCMD
). Scholars in economics, the life sciences, and many other fields tend to co-author disproportionately with same-gender colleagues (\mbox{%DIFAUXCMD
\citeproc{ref-mcdowell1992effect}{15}}\hskip0pt%DIFAUXCMD
}\DIFaddend --\DIFdelbegin \DIFdel{\mbox{%DIFAUXCMD
\citeproc{ref-yang2022gender}{14}}\hskip0pt%DIFAUXCMD
) , less is known about their teaching choices, including the formation of co-teaching teams , interdisciplinarity, and reading selections, and their effects on student outcomes. }\DIFdelend \DIFaddbegin \DIFadd{\mbox{%DIFAUXCMD
\citeproc{ref-kwiek2021gender}{18}}\hskip0pt%DIFAUXCMD
). Notably, this form of gender-based sorting is not merely driven by institutional constraints: it appears even in large and diverse research environments (e.g., \mbox{%DIFAUXCMD
\citeproc{ref-kwiek2021gender}{18}}\hskip0pt%DIFAUXCMD
) and across career stages (e.g., \mbox{%DIFAUXCMD
\citeproc{ref-nguyen2025mind}{19}}\hskip0pt%DIFAUXCMD
). Such gender homophily is also considered a contributing factor to gender gaps in outputs (e.g., \mbox{%DIFAUXCMD
\citeproc{ref-ductor2023gender}{20}}\hskip0pt%DIFAUXCMD
) and to broader gender inequalities in academia (e.g., \mbox{%DIFAUXCMD
\citeproc{ref-mcdowell1992effect}{15}}\hskip0pt%DIFAUXCMD
).
}\DIFaddend 

\DIFdelbegin \DIFdel{Our findings contribute to the literature on gender dynamics in academia, particularly research team formation (\mbox{%DIFAUXCMD
\citeproc{ref-milojevic2014principles}{13}}\hskip0pt%DIFAUXCMD
, \mbox{%DIFAUXCMD
\citeproc{ref-boschini2007team}{15}}\hskip0pt%DIFAUXCMD
, \mbox{%DIFAUXCMD
\citeproc{ref-love2022impact}{16}}\hskip0pt%DIFAUXCMD
), which has shown significant gender homophily--- a tendency to collaborate with colleagues of the same gender (\mbox{%DIFAUXCMD
\citeproc{ref-kwiek2021gender}{17}}\hskip0pt%DIFAUXCMD
).
It also }\DIFdelend \DIFaddbegin \DIFadd{While research collaborations have been widely investigated, collaboration patterns in teaching remain less well understood. On the one hand, homophily may shape teaching teams by reflecting the same individual preferences observed in research collaborations. On the other hand, teaching has distinctive features: its predominantly local and in-person nature could either constrain or amplify gender-based homophily. These considerations motivate our effort to explore the prevalence and gender composition of teaching teams at scale.
}

\subsection{\DIFadd{Representation of Women in Citations and Course Materials}}\label{representation-of-women-in-citations-and-course-materials}

\DIFadd{A substantial body of research }\DIFaddend examines gender differences in citation patterns \DIFdelbegin \DIFdel{, a }\DIFdelend \DIFaddbegin \DIFadd{(\mbox{%DIFAUXCMD
\citeproc{ref-ductor2023gender}{20}}\hskip0pt%DIFAUXCMD
--\mbox{%DIFAUXCMD
\citeproc{ref-card2022gender}{23}}\hskip0pt%DIFAUXCMD
), }\DIFaddend driving factor of the persistent gender bias in academia (\DIFdelbegin \DIFdel{\mbox{%DIFAUXCMD
\citeproc{ref-boivin2024sexism}{18}}\hskip0pt%DIFAUXCMD
), including tenure promotion (\mbox{%DIFAUXCMD
\citeproc{ref-weisshaar2017publish}{19}}\hskip0pt%DIFAUXCMD
), grant success (\mbox{%DIFAUXCMD
\citeproc{ref-bedi2012gender}{20}}\hskip0pt%DIFAUXCMD
), co-authorships (\mbox{%DIFAUXCMD
\citeproc{ref-boschini2007team}{15}}\hskip0pt%DIFAUXCMD
) , and peer recognition (\mbox{%DIFAUXCMD
\citeproc{ref-sarsons2017recognition}{21}}\hskip0pt%DIFAUXCMD
, \mbox{%DIFAUXCMD
\citeproc{ref-card2022gender}{22}}\hskip0pt%DIFAUXCMD
).
While prior research has primarily focused on these areas, we shift our attention to the citation gap in teaching. This issue may not only reinforce existing genderbias in academia but also shape students' learning outcomes and future career choices (\mbox{%DIFAUXCMD
\citeproc{ref-card2022impact}{23}}\hskip0pt%DIFAUXCMD
), effects that are less well understood. }\DIFdelend \DIFaddbegin \DIFadd{\mbox{%DIFAUXCMD
\citeproc{ref-boivin2024sexism}{24}}\hskip0pt%DIFAUXCMD
). Citation practices often reflect gender homophily (e.g., \mbox{%DIFAUXCMD
\citeproc{ref-card2022gender}{23}}\hskip0pt%DIFAUXCMD
), and male-dominated fields may confer a ``first-mover advantage'' on men (e.g., \mbox{%DIFAUXCMD
\citeproc{ref-newman2009first}{25}}\hskip0pt%DIFAUXCMD
). Other contributing factors include different self-citation patterns (\mbox{%DIFAUXCMD
\citeproc{ref-mishra2018self}{22}}\hskip0pt%DIFAUXCMD
) and overall gaps in publication output (\mbox{%DIFAUXCMD
\citeproc{ref-ductor2023gender}{20}}\hskip0pt%DIFAUXCMD
).
}\DIFaddend 

\DIFdelbegin \DIFdel{By focusing on the relationships between teaching and interdisciplinarity, our work also contributes to the literature on the factors influencing interdisciplinarity in academia. There is evidence that interdisciplinarity research is less likely to be funded (\mbox{%DIFAUXCMD
\citeproc{ref-bromham2016interdisciplinary}{24}}\hskip0pt%DIFAUXCMD
), tends to attract fewer citations when it is highly interdisciplinary (\mbox{%DIFAUXCMD
\citeproc{ref-lariviere2010relationship}{25}}\hskip0pt%DIFAUXCMD
), and is correlated with the probability of publication in academic journals (\mbox{%DIFAUXCMD
\citeproc{ref-xiang2025evaluating}{\textbf{xiang2025evaluating?}}}\hskip0pt%DIFAUXCMD
). Furthermore, students attending colleges }\DIFdelend \DIFaddbegin \DIFadd{Evidence that these disparities extend into teaching is more limited, but studies indicate that female scholars are underrepresented in course syllabi (\mbox{%DIFAUXCMD
\citeproc{ref-harris2020diversify}{26}}\hskip0pt%DIFAUXCMD
--\mbox{%DIFAUXCMD
\citeproc{ref-skitka2021gender}{28}}\hskip0pt%DIFAUXCMD
). This underrepresentation appears linked to the instructor's gender. For example, in Political Science, instructors' gender predicts the assignment of female-authored readings, suggesting a form of gender homophily in course content (\mbox{%DIFAUXCMD
\citeproc{ref-smith2020gender}{29}}\hskip0pt%DIFAUXCMD
). Similarly, in research on political science, (\mbox{%DIFAUXCMD
\citeproc{ref-dion2018gendered}{21}}\hskip0pt%DIFAUXCMD
) finds that female-led teams cite more women than male-led teams, whereas mixed-gender teams do not cite more women than all-male teams. These findings suggest that team composition and dynamics strongly influence decisions about what to cite or include in course readings. A key goal of our study is to investigate these patterns on a large scale across multiple domains and institutions.
}

\subsection{\DIFadd{Mixed-Gender Teams and Interdisciplinarity}}\label{mixed-gender-teams-and-interdisciplinarity}

\DIFadd{Another relevant strand of literature concerns the drivers of interdisciplinary work in research. Recent evidence suggests that women, on average, undertake more interdisciplinary research than men (\mbox{%DIFAUXCMD
\citeproc{ref-pinheiro2022women}{30}}\hskip0pt%DIFAUXCMD
), and that interdisciplinary collaborations can produce broader scientific outcomes (\mbox{%DIFAUXCMD
\citeproc{ref-specht2022interdisciplinary}{31}}\hskip0pt%DIFAUXCMD
). Taken together, this work suggests that the presence of women in a team is associated }\DIFaddend with more interdisciplinary \DIFdelbegin \DIFdel{courses tend to earn higher earnings after graduation (\mbox{%DIFAUXCMD
\citeproc{ref-han2023interdisciplinary}{3}}\hskip0pt%DIFAUXCMD
, \mbox{%DIFAUXCMD
\citeproc{ref-evans2016measuring}{26}}\hskip0pt%DIFAUXCMD
).
We extend this work by examining the association between gender composition and interdisciplinarity in teaching, showing that mixed-gender teams occur less frequently but tend to be more interdisciplinary . This finding suggests that removing barriers to mixed-gender team formation may increase interdisciplinarity}\DIFdelend \DIFaddbegin \DIFadd{intellectual contributions.
}

\DIFadd{In a teaching context, however, other factors may influence this relationship. For example, being interdisciplinary in research does not necessarily translate into interdisciplinary teaching styles. Furthermore, students may not demand interdisciplinary courses. Or there might be institutional constraints that prevent programs from supporting interdisciplinarity. Finally, team dynamics or matching patterns may prevent instructors from finding good partners for interdisciplinary courses. To shed light on these issues, we explored, at a large scale, the relationships among course interdisciplinarity, team size, and gender composition }\DIFaddend in teaching.

\DIFdelbegin \DIFdel{We also build on prior research on the underrepresentation of female-authored works in university curricula, which has been shown within specific fields (psychology and international relations)and at a small scale (\mbox{%DIFAUXCMD
\citeproc{ref-skitka2021gender}{27}}\hskip0pt%DIFAUXCMD
}\DIFdelend \DIFaddbegin \subsection{\DIFadd{Gender Diversity and Innovation}}\label{gender-diversity-and-innovation}

\DIFadd{Extensive research has shown that gender-diverse teams tend to generate more novel and higher-impact ideas than single-gender teams in scientific collaboration (e.g., \mbox{%DIFAUXCMD
\citeproc{ref-yang2022gender}{32}}\hskip0pt%DIFAUXCMD
). Related work on research and R\&D teams links gender diversity to more radical or innovative outputs, as well as to improved decision-making and performance (\mbox{%DIFAUXCMD
\citeproc{ref-apesteguia2012impact}{33}}\hskip0pt%DIFAUXCMD
}\DIFaddend --\DIFdelbegin \DIFdel{\mbox{%DIFAUXCMD
\citeproc{ref-harris2020diversify}{29}}\hskip0pt%DIFAUXCMD
).
Our results reveal a consistent and significant gap in the number of cited works between female and male instructors, which is only partially addressed within mixed-gender co-teaching. This evidence highlights teaching practices as an important factor shaping gender differences in recognition among scholars. }\DIFdelend \DIFaddbegin \DIFadd{\mbox{%DIFAUXCMD
\citeproc{ref-love2022impact}{37}}\hskip0pt%DIFAUXCMD
).
}\DIFaddend 

\DIFdelbegin \DIFdel{Finally, studies about team formation are especially relevant to our work. These studies have found that }\DIFdelend \DIFaddbegin \DIFadd{However, in the context of teaching, the relationship between gender diversity and ``novelty'' is less well understood. A key challenge lies in defining measures of novelty in a pedagogical setting. First, syllabi often rely heavily on textbooks rather than on journal articles, and pedagogical innovation may lag behind the scientific frontier. Moreover, innovation in teaching can take different forms: instructors can adopt new methods that are pedagogically novel but not at the scientific frontier (``teaching novelty''); they can assign more recent materials (``recency''); or they can include papers at the forefront of research (``scientific novelty''). In this study, we propose simple measures to begin empirically exploring these three dimensions at scale.
}

\subsection{\DIFadd{Covid-19's impact on academic work}}\label{covid-19s-impact-on-academic-work}

\DIFadd{COVID-19 disrupted academic work, reducing research time and scientific output, with particularly strong effects for women and researchers with caregiving responsibilities (\mbox{%DIFAUXCMD
\citeproc{ref-myers2020unequal}{11}}\hskip0pt%DIFAUXCMD
, \mbox{%DIFAUXCMD
\citeproc{ref-viglione2020women}{38}}\hskip0pt%DIFAUXCMD
--\mbox{%DIFAUXCMD
\citeproc{ref-garcia2024silver}{41}}\hskip0pt%DIFAUXCMD
). At the same time, the pandemic also reshaped collaboration networks and scientific meetings: studies report a decrease in new co-authorships (\mbox{%DIFAUXCMD
\citeproc{ref-gao2021potentially}{39}}\hskip0pt%DIFAUXCMD
) and an increase in cross-border teaming (\mbox{%DIFAUXCMD
\citeproc{ref-venturini2025effectsremoteworkingscientific}{42}}\hskip0pt%DIFAUXCMD
), and a shift towards virutal conferences and online meetings. While online formats broadened participation, with a modest increase in women's representation (\mbox{%DIFAUXCMD
\citeproc{ref-olechnicka2025impact}{43}}\hskip0pt%DIFAUXCMD
), they also reduced opportunities for informal, in-person interactions that are seen as vital for collaboration and idea generation (\mbox{%DIFAUXCMD
\citeproc{ref-bray2022together}{44}}\hskip0pt%DIFAUXCMD
). In parallel, universities were forced into a rapid transition to online and hybrid teaching, potentially altering how teaching teams are formed and how often instructors collaborate across gender and disciplinary boundaries. We therefore exploit the onset of COVID-19 as an exogenous shock and estimate trends around 2020 to examine whether the prevalence and composition of }\DIFaddend mixed-gender \DIFdelbegin \DIFdel{teams tend to perform better in various settings (\mbox{%DIFAUXCMD
\citeproc{ref-yang2022gender}{14}}\hskip0pt%DIFAUXCMD
, \mbox{%DIFAUXCMD
\citeproc{ref-bear2011role}{30}}\hskip0pt%DIFAUXCMD
)}\DIFdelend \DIFaddbegin \DIFadd{teams changed after the pandemic, conditional on institution and field}\DIFaddend .
\DIFdelbegin \DIFdel{In research collaboration, for example, mixed-gender teams often obtain more citations, produce more novel research , publish in more prestigious journals, and are more interdisciplinary (\mbox{%DIFAUXCMD
\citeproc{ref-yang2022gender}{14}}\hskip0pt%DIFAUXCMD
, \mbox{%DIFAUXCMD
\citeproc{ref-maddi2021gender}{31}}\hskip0pt%DIFAUXCMD
). Although we do not measure team ``performance'' as we lack data on students outcomes, our results show that mixed-gender co-teaching teams are consistently less likely to form, but, once created, they tend to deliver different outcomes, especially a higher interdisciplinarity, that previous literature has suggested having an impact on students outcomes as discussed above.
}\DIFdelend 

\section{Materials and Methods}\label{sec:methods}

\subsection{Data \DIFaddbegin \DIFadd{sources}\DIFaddend }\DIFdelbegin %DIFDELCMD < \label{data}
%DIFDELCMD < %%%
\DIFdelend \DIFaddbegin \label{data-sources}
\DIFaddend 

\DIFdelbegin \DIFdel{We obtained }\DIFdelend \DIFaddbegin \subsubsection{\DIFadd{Open Syllabus corpus}}\label{open-syllabus-corpus}

\DIFadd{The core dataset consists of }\DIFaddend a corpus of \DIFdelbegin \DIFdel{over six million documents }\DIFdelend \DIFaddbegin \DIFadd{more than six million university syllabi }\DIFaddend compiled by Open Syllabus (New York, US) \DIFdelbegin \DIFdel{. Open Syllabus created this dataset }\DIFdelend through web extractions \DIFdelbegin \DIFdel{that identified syllabi }\DIFdelend from university websites\DIFdelbegin \DIFdel{, }\DIFdelend \DIFaddbegin \DIFadd{. These extractions identify syllabi }\DIFaddend with a median confidence level of 99.8\%. A tagging algorithm \DIFdelbegin \DIFdel{extracted }\DIFdelend \DIFaddbegin \DIFadd{extracts }\DIFaddend key course details, such as the title, \DIFdelbegin \DIFdel{field, }\DIFdelend description, academic year, \DIFdelbegin \DIFdel{duration, and }\DIFdelend language, a list of anonymised instructors, and the assigned readings.\footnote{\DIFdelbegin \DIFdel{The documentation }\DIFdelend \DIFaddbegin \DIFadd{Documentation }\DIFaddend available at: \url{https://docs.opensyllabus.org}} \DIFaddbegin \DIFadd{OpenSyllabus also assigns each syllabus to one of 69 top-level disciplinary fields, derived from the U.S. Department of Education's CIP classification system.}\footnote{
  \DIFadd{The CIP classification is available at:
  }\url{https://nces.ed.gov/ipeds/cipcode/browse.aspx?y=55}}

\DIFaddend While the original dataset \DIFdelbegin \DIFdel{included }\DIFdelend \DIFaddbegin \DIFadd{includes }\DIFaddend syllabi in 49 languages, most documents (96\%) \DIFdelbegin \DIFdel{were }\DIFdelend \DIFaddbegin \DIFadd{are }\DIFaddend in English. For simplicity, we focused exclusively on \DIFdelbegin \DIFdel{these documents.}\DIFdelend \DIFaddbegin \DIFadd{English-language syllabi.}\DIFaddend \footnote{
  Our focus on English-language courses means that, while we have comprehensive data for English-speaking countries such as Canada, Ireland, the United States, and Great Britain, the sample in non-English-speaking countries tends to be more representative of internationally oriented universities. These are typically institutions offering programs in English or advanced-level courses, such as postgraduate programs or disciplines where English is the primary medium of instruction within traditional universities.} \DIFdelbegin %DIFDELCMD < 

%DIFDELCMD < %%%
\DIFdel{The resulting dataset comprised }\DIFdelend \DIFaddbegin \DIFadd{This choice resulted in a total of }\DIFaddend 5.4 million syllabi \DIFaddbegin \DIFadd{(1990-2019) }\DIFaddend from approximately 4,000 higher education institutions across fifteen countries (\DIFdelbegin \DIFdel{see Table \ref{tab:SI-country}for the list of countries)from 1990 to 2019, with a higher proportion of syllabi in recent years (see Table \ref{tab:SI-years}).OpenSyllabus classified these syllabi into 69 top-level fields derived from the U.S. Department of Education's CIP code classification (see Table \ref{tab:SI-fields}) }\DIFdelend \DIFaddbegin \DIFadd{Table \ref{tab:SI-country})}\DIFaddend .
\DIFdelbegin \footnote{\DIFdel{The original CIP classification is available online: }%DIFDELCMD < \url{https://nces.ed.gov/ipeds/cipcode/browse.aspx?y=55}%%%
} %DIFAUXCMD
\addtocounter{footnote}{-1}%DIFAUXCMD
\DIFdel{About 2.9 million syllabi (53\% of the total) listed readings matched with bibliographic sources, providing additional metadata about authorship information, journal, and publication year. The institution was matched to a list of more than 22, 000 entities from the Research Organisation Registry, providing further metadata, including the institution's country and enrollment figures --- the institutions in our sample account for over 35 million enrolled students today.}%DIFDELCMD < \footnote{%%%
\DIFdel{For }\DIFdelend \DIFaddbegin 

\DIFadd{Because course level (e.g., undergraduate or graduate) was not explicitly reported in the data, we inferred it from course codes. For example, ``Psychology 101'' typically denotes an introductory course, whereas ``Psychology 401'' indicates a }\DIFaddend more \DIFdelbegin \DIFdel{details on how OpenSyllabus classified and matched the data}\DIFdelend \DIFaddbegin \DIFadd{advanced offering. Using this heuristic, we categorized syllabi into four levels: basic, advanced, graduate, or unknown (see Section \ref{course-level} for the details).}\footnote{
  \DIFadd{According to (\mbox{%DIFAUXCMD
\citeproc{ref-biasi2022education}{10}}\hskip0pt%DIFAUXCMD
)}\DIFaddend , \DIFdelbegin \DIFdel{the related documentation }\DIFdelend \DIFaddbegin \DIFadd{this classification procedure }\DIFaddend is \DIFdelbegin \DIFdel{available online at }%DIFDELCMD < \url{https://docs.opensyllabus.org}%%%
\DIFdelend \DIFaddbegin \DIFadd{approximately 85\% accurate for a sample of U}\DIFaddend .\DIFaddbegin \DIFadd{S. courses that could be matched to catalogues with verifiable information.}\DIFaddend }

\DIFaddbegin \subsubsection{\DIFadd{Instructor genders}}\label{instructor-genders}

\DIFaddend Each syllabus lists one or more instructors, with 76\% listing a single instructor, 16\% listing two, 4\% listing three, and 4\% listing more than three instructors. \DIFdelbegin \DIFdel{OpenSyllabus determined the instructorsgender automatically }\DIFdelend \DIFaddbegin \DIFadd{For each syllabus, OpenSyllabus automatically determines the instructors' gender }\DIFaddend based on their names\DIFdelbegin \DIFdel{, resulting in }\DIFdelend \DIFaddbegin \DIFadd{. This gender disambiguation process employs state-of-the-art methods that have been shown to achieve high accuracy on other datasets (e.g., \mbox{%DIFAUXCMD
\citeproc{ref-dion2018gendered}{21}}\hskip0pt%DIFAUXCMD
), which are detailed in the Supplementary Information (Section \ref{gender-disambiguation}). In our dataset, this classification resulted in }\DIFaddend 52\% male, 37\% female, and 11\% unknown categories. (Table \ref{tab:SI-composition} illustrates the gender composition of the teaching teams). After excluding syllabi with unknown gender, the distribution was 58\% male and 42\% female instructors\DIFdelbegin \DIFdel{, which aligns closely }\DIFdelend \DIFaddbegin \DIFadd{.
}

\DIFadd{These figures broadly align }\DIFaddend with the 45\% of female academic staff reported in OECD countries (\DIFdelbegin \DIFdel{\mbox{%DIFAUXCMD
\citeproc{ref-oecd2022}{32}}\hskip0pt%DIFAUXCMD
). Open Syllabus used the same method to determine the gender of the authors listed in the readings}\DIFdelend \DIFaddbegin \DIFadd{\mbox{%DIFAUXCMD
\citeproc{ref-oecd2022}{45}}\hskip0pt%DIFAUXCMD
). Although official statistics to assess the representativeness of this gender population are not available, a detailed comparison with official statistics (Supporting Info, Section \ref{representativeness-of-the-gender-distribution}) shows that the observed trends in the proportions of women and men are highly correlated with representative statistics on academic staff across countries and with the distribution of graduates across fields (used as a proxy for the number of instructors). This analysis indicates a reasonable representativeness of our data.
}

\subsubsection{\DIFadd{Linking syllabi to bibliometric data}}\label{linking-syllabi-to-bibliometric-data}

\DIFadd{About 2.9 million syllabi (53\% of the total) list readings matched with bibliographic sources, providing additional metadata on authorship, journal of publication, and publication year. Institutions are matched to a list of more than 22}\DIFaddend ,\DIFdelbegin \DIFdel{resulting in 32\% female and 56\% male authors, with 12\% of the authors having an unknown gender}\DIFdelend \DIFaddbegin \DIFadd{000 entities from the Research Organisation Registry (ROR), providing further metadata on institutional characteristics, including country and enrollment figures. Overall, the institutions in our sample account for over 35 million enrolled students today.}\footnote{
  \DIFadd{For more details on how OpenSyllabus classify and match the data, the related documentation is available online at }\url{https://docs.opensyllabus.org}\DIFadd{.}}\DIFadd{. Overall, our dataset captures roughly 50\% of the educational institutions listed in ROR, underscoring its broad coverage.}\footnote{
  \DIFadd{The same methods used to identify the gender of the instructors were used for the gender of the authors listed in the readings, resulting in 32\% female and 56\% male authors, with 12\% of the authors having an unknown gender.}}

\DIFadd{To assess the scientific novelty of the readings assigned in each syllabus, we further linked the data to SciSciNet (v2) (\mbox{%DIFAUXCMD
\citeproc{ref-lin2023sciscinet}{46}}\hskip0pt%DIFAUXCMD
), a large-scale database of scientific publications that includes pre-computed indicators of novelty, such as the atypicality measure introduced by Uzzi et al. (\mbox{%DIFAUXCMD
\citeproc{ref-uzzi2013atypical}{47}}\hskip0pt%DIFAUXCMD
). Matches were performed using article DOIs. Unfortunately, most syllabi include textbooks that do not have a DOI and a corresponding atypicality measure. Even so, the process resulted in a match with 80\% of all readings with at least an available DOI in the readings. These matched publications appear in roughly 400,000 syllabi, representing about 14\% of the 2.9 million syllabi containing reading lists.
}

\subsubsection{\DIFadd{Supplementary Covid-19 dataset}}\label{supplementary-covid-19-dataset}

\DIFadd{To examine changes associated with the Covid-19 lockdown period, we supplement this dataset with a second corpus consisting of 2.2 million syllabi of courses between 2019 and 2023 from 33 European countries (EU-27 plus Switzerland, Norway, Iceland, the United Kingdom, and Liechtenstein)}\DIFaddend .

\subsection{Outcome variables}\label{outcome-variables}

Building on prior research highlighting the value of \DIFdelbegin \DIFdel{teaching }\DIFdelend \DIFaddbegin \DIFadd{exposing student to }\DIFaddend cutting-edge knowledge (\citeproc{ref-biasi2022education}{10}), the \DIFdelbegin \DIFdel{role }\DIFdelend \DIFaddbegin \DIFadd{influence }\DIFaddend of interdisciplinary training \DIFdelbegin \DIFdel{in shaping }\DIFdelend \DIFaddbegin \DIFadd{on }\DIFaddend career outcomes (\citeproc{ref-han2023interdisciplinary}{3}), and evidence of gender disparities in \DIFdelbegin \DIFdel{citation practices---particularly the under-recognition of female scholars and the influence of the citing author's gender (\mbox{%DIFAUXCMD
\citeproc{ref-aksnes2011female}{33}}\hskip0pt%DIFAUXCMD
}\DIFdelend \DIFaddbegin \DIFadd{academic recognition (\mbox{%DIFAUXCMD
\citeproc{ref-aksnes2011female}{48}}\hskip0pt%DIFAUXCMD
}\DIFaddend --\DIFdelbegin \DIFdel{\mbox{%DIFAUXCMD
\citeproc{ref-tekles2022same}{35}}\hskip0pt%DIFAUXCMD
)--- we defined key outcome variables to }\DIFdelend \DIFaddbegin \DIFadd{\mbox{%DIFAUXCMD
\citeproc{ref-tekles2022same}{50}}\hskip0pt%DIFAUXCMD
), we focus on a set of outcome variables designed to capture these dimensions of the assigned readings and course content. These variables allow us to }\DIFaddend examine how team configurations, including gender composition and team size, are associated with the interdisciplinarity of assigned readings and their key characteristics, such as publication age and the inclusion of female-authored works. The outcome variables are described in detail below; a concise summary is \DIFdelbegin \DIFdel{provided in Table \ref{tab:outcome-variables}}\DIFdelend \DIFaddbegin \DIFadd{in Table \ref{tab:outcomes}}\DIFaddend .

\begin{table}
\centering\centering
\caption{\DIFdelbeginFL %DIFDELCMD < \label{tab:outcome-variables}%%%
\DIFdelendFL \DIFaddbeginFL \label{tab:outcomes}\DIFaddendFL Outcome variables}
\centering
\DIFdelbeginFL %DIFDELCMD < \fontsize{12}{14}\selectfont
%DIFDELCMD < \begin{tabular}[t]{>{\raggedright\arraybackslash}p{4cm}>{\raggedright\arraybackslash}p{10cm}}
%DIFDELCMD < \toprule
%DIFDELCMD < \multicolumn{1}{>{}p{4cm}}{\textbf{Variable}} & \multicolumn{1}{>{}p{10cm}}{\textbf{Definition}}\\
%DIFDELCMD < \midrule
%DIFDELCMD < \textbf{Interdisciplinarity} & Percentile rank of the course's interdisciplinarity score for the year.\\
%DIFDELCMD < \addlinespace
%DIFDELCMD < \textbf{Age of Readings} & Percentile rank of the average publication age of assigned readings.\\
%DIFDELCMD < \addlinespace
%DIFDELCMD < \textbf{Ratio of Female Authors} & Proportion of women authors in the assigned readings.\\
%DIFDELCMD < \bottomrule
%DIFDELCMD < \end{tabular}
%DIFDELCMD < %%%
\DIFdelendFL \DIFaddbeginFL \begin{tabu} to \linewidth {>{\raggedright\arraybackslash}p{4cm}>{\raggedright\arraybackslash}p{10cm}}
\toprule
\multicolumn{1}{>{}p{4cm}}{\textbf{Variable}} & \multicolumn{1}{>{}p{10cm}}{\textbf{Definition}}\\
\midrule
\textbf{Course Interdisciplinarity} & Percentile rank of the course's interdisciplinarity score for the year.\\
\addlinespace
\textbf{Age of Readings} & Percentile rank of the average publication age of assigned readings.\\
\addlinespace
\textbf{Atypicality of Readings} & Percentile rank of the median atypicality of assigned readings\\
\addlinespace
\textbf{Conventionality of Readings} & Percentile rank of the median citation of the assigned readings as referenced in other syllabi\\
\addlinespace
\textbf{Ratio of Female Authors} & Proportion of women authors in the assigned readings.\\
\bottomrule
\end{tabu}
\DIFaddendFL 
\end{table}

\subsubsection{Interdisciplinarity}\label{interdisciplinarity}

We measured interdisciplinarity by analysing course descriptions from syllabi, following Evans et al. (\citeproc{ref-han2023interdisciplinary}{3}, \DIFdelbegin \DIFdel{\mbox{%DIFAUXCMD
\citeproc{ref-evans2016measuring}{26}}\hskip0pt%DIFAUXCMD
}\DIFdelend \DIFaddbegin \DIFadd{\mbox{%DIFAUXCMD
\citeproc{ref-evans2016measuring}{51}}\hskip0pt%DIFAUXCMD
}\DIFaddend ). This approach converts course descriptions into ``bags of words'', where word frequencies are normalised by the inverse ratio of term frequency to document frequency (TF-IDF). We then computed a correlation matrix of word usage across academic fields to quantify the conceptual distances between them. The interdisciplinarity score for each syllabus was calculated as a weighted average of pairwise correlations with other syllabi, with weights reflecting field distance (see \DIFdelbegin \DIFdel{Supporting Information, }\DIFdelend Section \ref{si-interdisciplinarity}). This method ensures that syllabi associated with distant fields, either academically or conceptually, are considered more interdisciplinary.\footnote{
  To scale to millions of documents, we used random subsampling of academic fields over years and averaged results across subsamples for robustness\DIFdelbegin \DIFdel{(see Supporting Information}\DIFdelend , \DIFaddbegin \DIFadd{as outlined in }\DIFaddend Section \ref{si-interdisciplinarity}\DIFdelbegin \DIFdel{)}\DIFdelend .}

We expressed interdisciplinarity for each syllabus \(i\) as the percentile rank within each academic year of the interdisciplinarity score:
\[
  \text{Interdisciplinarity}_i = \text{PR}_{yr} (\text{Interdisciplinarity Score}_i),
\]
Where \(\text{PR}_{yr}\) represents the percentile rank function applied to all syllabi within a given year \(yr\). This normalisation controls for skewed score distributions and allows comparison across years.

\subsubsection{\DIFaddbegin \DIFadd{Recency, Conventionality, and Atypicality of the }\DIFaddend Readings\DIFdelbegin \DIFdel{Selection}\DIFdelend }\DIFdelbegin %DIFDELCMD < \label{readings-selection}
%DIFDELCMD < %%%
\DIFdelend \DIFaddbegin \label{recency-conventionality-and-atypicality-of-the-readings}
\DIFaddend 

To characterise the \DIFaddbegin \DIFadd{novelty of the }\DIFaddend assigned readings, we \DIFdelbegin \DIFdel{calculated two }\DIFdelend \DIFaddbegin \DIFadd{used three }\DIFaddend key dimensions: \DIFdelbegin \DIFdel{the age of the readings and the proportion of female authors represented.
}\DIFdelend \DIFaddbegin \DIFadd{(1) the publication year of the assigned readings (2) the frequency of new articles or books introduced in the syllabus; and (3) the median scientific novelty of the readings as measured by the Uzzi et al's atypicality measure (\mbox{%DIFAUXCMD
\citeproc{ref-uzzi2013atypical}{47}}\hskip0pt%DIFAUXCMD
).
}

\DIFaddend First, we define the total number of references, \(N_i\), as the sum of articles, books, and chapters listed in a syllabus \(i\). This variable serves as a broad measure of a course's ``breadth,'' as more assigned readings may indicate a more extensive or comprehensive curriculum. Then, we define the \emph{Age of Readings} variable as the difference between the syllabus year (\(\text{Year}_i\)) and the average publication year of the assigned readings \(k=1,\cdots,K\):\\
\[
  \text{Age of Readings}_i = \text{PR}_{yr}\left(\text{Year}_i - \sum_{k=1}^{N_i} \text{Publication Year}_k / N_i\right),
\]
Where \(\text{PR}_{yr}\) represents the percentile rank function applied to all syllabi within a given year \(yr\). This variable provides a proxy of how recent, or ``novel,'' the readings are.
\DIFdelbegin \DIFdel{While more sophisticated methods to measure novelty are available (\mbox{%DIFAUXCMD
\citeproc{ref-uzzi2013atypical}{36}}\hskip0pt%DIFAUXCMD
, \mbox{%DIFAUXCMD
\citeproc{ref-wang2017bias}{37}}\hskip0pt%DIFAUXCMD
), we employ this more straightforward approach because pedagogical innovation tends to be incremental, and computationally intensive methods are impractical for large-scale datasets like ours. We anticipate that novelty will vary across fields (e. g., history courses may cite older texts more frequently than computer science courses). Still, within-field differences in novelty serve as a proxy for how close a course is to the current frontier of knowledge (``novelty'')}\DIFdelend \DIFaddbegin 

\DIFadd{Second, to capture the conventionality of the readings, we compute the }\emph{\DIFadd{Share of Conventional Readings}} \DIFadd{as the proportion of references in syllabus \(i\) that appear in any earlier iteration of the same course (as observed in our dataset). This measure reflects how frequently instructors refresh their syllabi by incorporating newly selected books or articles, which is an important dimension of curricular innovation that is not captured by publication dates alone.
}

\DIFadd{Third, to quantify the scientific novelty of the assigned readings, we use the median atypicality score of all references in syllabus \(i\). Atypicality, introduced by Uzzi et al. (\mbox{%DIFAUXCMD
\citeproc{ref-uzzi2013atypical}{47}}\hskip0pt%DIFAUXCMD
), measures how unexpectedly different a paper's combination of prior knowledge is relative to conventional citation patterns. We rely on the pre-computed atypicality values provided in SciSciNet (v2) by Lin et al. (\mbox{%DIFAUXCMD
\citeproc{ref-lin2023sciscinet}{46}}\hskip0pt%DIFAUXCMD
). Higher values indicate that the readings draw from more unconventional or cross-disciplinary combinations of ideas, offering a complementary dimension of novelty that is orthogonal to recency or syllabus updating.}\footnote{
  \DIFadd{Note that this measure only captures the atypicality at the time of publication, hence, it may not entirely reflect the atypicality at the time of the course. This issue can be partly accounted for in the analysis by (1) aggregating over multiple readings and (2) controlling for the average publication year of the readings or alternatively by imposing a discount for articles published beyond a certain threshold from the course date (e.g., imposing a mean atypicality for readings five years or more older than the course year).}}

\DIFadd{Together, these three measures, recency of publications, the conventionality of readings, and the scientific atypicality of the readings, provide a multidimensional characterisation of curricular novelty}\DIFaddend .

\DIFaddbegin \subsubsection{\DIFadd{Women Authors Cited}}\label{women-authors-cited}

\DIFaddend Finally, we define the \emph{Ratio of Female Authors} as the proportion of female authors among all authors in the assigned readings:
\[
  \text{Ratio of Female Authors}_i = \frac{\text{Female Authors}_i + 1}{\text{Female Authors}_i + \text{Male Authors}_i + 2}.
\]
Here, we add two pseudo-observations (one for each gender) to stabilise the ratio, preventing extreme values in cases with very few authors (\DIFdelbegin \DIFdel{\mbox{%DIFAUXCMD
\citeproc{ref-agresti2000simple}{38}}\hskip0pt%DIFAUXCMD
}\DIFdelend \DIFaddbegin \DIFadd{\mbox{%DIFAUXCMD
\citeproc{ref-agresti2000simple}{52}}\hskip0pt%DIFAUXCMD
}\DIFaddend ). This metric enables us to examine whether gender and collaboration are associated with the representation of female-authored work in teaching.

\subsection{Simulating Teaching Collaborations}\label{simulating-teaching-collaborations}

Building on previous research showing that gender-diverse teams are underrepresented in science (\DIFdelbegin \DIFdel{\mbox{%DIFAUXCMD
\citeproc{ref-guimera2005team}{11}}\hskip0pt%DIFAUXCMD
}\DIFdelend \DIFaddbegin \DIFadd{\mbox{%DIFAUXCMD
\citeproc{ref-guimera2005team}{53}}\hskip0pt%DIFAUXCMD
}\DIFaddend ), we investigate the evolution of the gender composition of co-teaching teams to determine whether it deviates from what is expected under a gender-neutral matching process. Crucially, this analysis accounts for the fact that team formation is influenced not only by individual preferences but also by institutional factors, such as departmental affiliation or academic field, as co-teaching typically occurs among faculty within the same department or the same academic field. To account for these institutional constraints, we employ a Monte Carlo randomisation approach adapted from Uzzi et al. (\DIFdelbegin \DIFdel{\mbox{%DIFAUXCMD
\citeproc{ref-guimera2005team}{11}}\hskip0pt%DIFAUXCMD
, \mbox{%DIFAUXCMD
\citeproc{ref-uzzi2013atypical}{36}}\hskip0pt%DIFAUXCMD
}\DIFdelend \DIFaddbegin \DIFadd{\mbox{%DIFAUXCMD
\citeproc{ref-uzzi2013atypical}{47}}\hskip0pt%DIFAUXCMD
, \mbox{%DIFAUXCMD
\citeproc{ref-guimera2005team}{53}}\hskip0pt%DIFAUXCMD
}\DIFaddend ). This approach preserves the overall gender composition and team size distribution while randomly reassigning instructors within strata defined by institution, field, and academic year. This simulation generates a randomised network in which instructors form teams independently of gender while maintaining the underlying institutional and disciplinary constraints. Therefore, by comparing actual data against teams formed in the randomised network, we can assess the difference between the observed team gender composition in the data and a counterfactual situation where instructors formed teams in a gender-neutral manner. See the details in the Supplementary Information (Section \ref{si-simulations}).

\subsection{Regression analysis}\label{regression-analysis}

To evaluate the association between gender composition, team size, and course content, we define four teaching team configurations, \(j\in \{F, M, MM, MF/FM, FF\}\), where F and M denote female and male instructors teaching alone, respectively. MM represents two male instructors; FM/MF indicates mixed-gender pairs; and FF corresponds to two female instructors. We then employ regression analysis to study the association between these team configurations and course outcomes across academic years, \DIFdelbegin \DIFdel{\(t = 1999, \cdots, 2020\)}\DIFdelend \DIFaddbegin \DIFadd{\(t = 2000, \cdots, 2019\)}\DIFaddend . As discussed before, the outcomes of interest include (1) interdisciplinarity, (2) the average age of readings, and (3) the proportion of cited female authors. To estimate these effects, we employ the following linear mixed-effects model:
\[
Y_{j, t} = 
 \alpha_t + 
 \text{Team}_{j, t} + 
 \text{STEM}_t + \eta_{f,t} +  % fields
 \text{Country}_{t} + 
 \DIFdelbegin \DIFdel{\text{Enroll}_t }\DIFdelend \DIFaddbegin \DIFadd{\text{CourseLevel}_{l,t} }\DIFaddend + 
 \delta_{t} + 
 \epsilon_{j, t}.
\]
Where:

\begin{itemize}
\item
  \(Y_{j, t, f, c, s}\) is the outcome variable for team configuration \(j\) in year \(t\), field \(f\), and school \(s\),
\item
  \DIFdelbegin \DIFdel{\(\alpha_t\) }\DIFdelend \DIFaddbegin \DIFadd{\(\alpha_{t}\) }\DIFaddend is a fixed effect for the academic year \(t\),
\item
  \(\text{Team}_{j, t}\) is a fixed effect for team configuration \(j\) in year \(t\),
\item
  \DIFdelbegin \DIFdel{\(STEM_t\) }\DIFdelend \DIFaddbegin \DIFadd{\(STEM_{t}\) }\DIFaddend is a fixed effect for STEM courses in year \(t\),
\item
  \(\eta_{f,t}\) is a random intercept for academic field \(f = 1,\cdots, 69\) at year \(t\),
\item
  \(\text{Country}_{c, t}\) is fixed effects for the institution \(s\)'s country \(c\) in year \(t\),
\item
  \DIFdelbegin \DIFdel{\(\text{Enroll}_{t}\) }\DIFdelend \DIFaddbegin \DIFadd{\(\text{CourseLevel}_{l, t}\) }\DIFaddend is an effect for the \DIFdelbegin \DIFdel{institution's enrollment size }\DIFdelend \DIFaddbegin \DIFadd{level of the course (\(l\) = basic, advanced, graduate, unknown) }\DIFaddend at year \(t\),
\item
  \(\delta_{s, t}\) is a random intercept for each of approximately 4000 unique institutions,
\item
  \(\epsilon_{j,t}\) is the residual error term.
\end{itemize}

Together, the model allows us to examine the association between teaching team configurations and citation-related course outcomes while accounting for potential confounding factors across disciplines and institutions.

\section{Results}\label{results}

\begin{figure}

{\centering \DIFdelbeginFL %DIFDELCMD < \includegraphics[width=1\linewidth]{/Users/andreablasco/Documents_local/Work/02_Projects/01_Education_Projects/03_OpenSyllabus/gender-teaching-prefs/output/figures//evolution-1} 
%DIFDELCMD < %%%
\DIFdelendFL \DIFaddbeginFL \includegraphics[width=1\linewidth]{/Users/andreablasco/Documents_local/Projects/01_education/03_OpenSyllabus/gender-teaching-prefs/output/figures//evolution-1} 
\DIFaddendFL 

}

\caption{Evolution of Teaching Teams by Gender Composition and Size. Trends in course gender composition from 2000 to 2019 reveal a steady increase in women's participation (both in solo-taught and two-instructor courses), while the proportion of mixed-gender teams has remained relatively stable over time. (A) Percentage of solo-taught courses by instructor gender over time. (B) Percentage of two-instructor courses by gender composition over time. N = 5.1 million courses. The shaded area represents 95\% confidence interval of the proportion per academic year.}\label{fig:team-trends}
\end{figure}

Figure \ref{fig:team-trends} illustrates gender composition and team size trends in university courses for single- and two-instructor courses (2000-2019). Panel A shows a decline in male-only single-instructor courses, from approximately 60\% in 2000 to 46\% in 2019, alongside an increase in female-only single-instructor courses, from roughly 25\% to 36\%.\footnote{This overall trend aligns with the increase in women holding academic positions worldwide. In 2000, women comprised approximately 35\% of the global academic staff. By 2022, this figure had risen to about 44\%, according to data from the World Bank. The data are available at: \url{https://data.worldbank.org/indicator/SE.TER.TCHR.FE.ZS}} Panel B shows male-only two-instructor courses also decreased, from about 8\% in the early 2000s to 6\% in 2020, while female-only and mixed-gender two-instructor courses remained relatively stable, each accounting for approximately 4\% and 3\%, respectively. These patterns indicate a consistent shift toward greater gender diversity in course instruction over time, with mixed-gender teaching teams (3\%) consistently less common than same-gender teams (10\%).

\begin{figure}

{\centering \DIFdelbeginFL %DIFDELCMD < \includegraphics[width=1\linewidth]{/Users/andreablasco/Documents_local/Work/02_Projects/01_Education_Projects/03_OpenSyllabus/gender-teaching-prefs/output/figures//montecarlo-1} 
%DIFDELCMD < %%%
\DIFdelendFL \DIFaddbeginFL \includegraphics[width=1\linewidth]{/Users/andreablasco/Documents_local/Projects/01_education/03_OpenSyllabus/gender-teaching-prefs/output/figures//montecarlo-1} 
\DIFaddendFL 

}

\caption{Comparison of gender composition between actual courses and courses simulated with Monte Carlo. Simulations ensure gender neutrality while keeping constant institutional constraints (i.e., number of teams by size, institution, academic year, and 69 academic fields). Simulated mixed-gender teams consistently exceed observed proportions, highlighting persistent barriers to forming gender-diverse teaching teams.}\label{fig:montecarlo}
\end{figure}

To assess whether gender imbalances in co-teaching teams persist after accounting for institutional constraints across fields and institutions, Figure \ref{fig:montecarlo} compares observed teaching team compositions (2000--2019) with those generated by a gender-neutral Monte Carlo simulation. This model randomly shuffles instructors while preserving the overall distribution of course loads per year by institution and academic field but ignoring gender-based preferences or constraints (see Section \ref{sec:methods}).

The results reveal that actual data systematically deviate from simulated expectations: same-gender teams, such as female-female (FF) and male-male (MM) courses, occur \emph{more} frequently in actual data than in simulations by 66\% and 20\%, respectively, in 2019. By contrast, mixed-gender collaborations (MF/FM) are consistently \emph{underrepresented} in the actual data (3\%) compared to simulations (6\%). These findings suggest that significant social or institutional dynamics---beyond chance---reinforce same-gender pairings, especially among women. Such dynamics may include gender-based homophily, mentoring networks, or within-department barriers or assignment practices.

\begin{figure}

{\centering \DIFdelbeginFL %DIFDELCMD < \includegraphics[width=1\linewidth]{/Users/andreablasco/Documents_local/Work/02_Projects/01_Education_Projects/03_OpenSyllabus/gender-teaching-prefs/output/figures//montecarlo-by-field-1} 
%DIFDELCMD < %%%
\DIFdelendFL \DIFaddbeginFL \includegraphics[width=1\linewidth]{/Users/andreablasco/Documents_local/Projects/01_education/03_OpenSyllabus/gender-teaching-prefs/output/figures//montecarlo-by-field-1} 
\DIFaddendFL 

}

\caption{This figure disaggregates the comparison of proportions between actual (green) and simulated (orange) mixed-gender (MF/FM) courses by academic field. The shaded bar indicates a 95\% confidence level of the proportion.}\label{fig:montecarlo-by-field}
\end{figure}

Figure \ref{fig:montecarlo-by-field} disaggregates the comparison between actual and simulated mixed-gender (MF/FM) courses by academic field. Nearly all fields exhibit a stark underrepresentation of mixed-gender teams (relative to expectations) by a factor of two to four times. For instance, in Medicine (Health and Welfare), the actual share of mixed-gender collaborations is 4.8\% compared to an expected 9.1\%, indicating a significant imbalance. Similar gaps are evident in Law (3.8\% actual vs.~6.4\% simulated) and Linguistics (4.3\% vs.~6.7\%), underscoring the pervasiveness of barriers to mixed-gender teams across various academic fields. Remarkably, we observe no differences in these patterns between male-dominated and female-dominated fields (e.g., Engineering or Accounting vs.~Nursing or Chemistry), suggesting that the underrepresentation of mixed-gender teams may be a structural feature rather than one driven by field-specific gender imbalances.

\begin{figure}
\centering
\DIFdelbeginFL %DIFDELCMD < \pandocbounded{\includegraphics[keepaspectratio]{/Users/andreablasco/Documents_local/Work/02_Projects/01_Education_Projects/03_OpenSyllabus/gender-teaching-prefs/output/figures//sim-by-cntry-1.pdf}}
%DIFDELCMD < %%%
\DIFdelendFL \DIFaddbeginFL \pandocbounded{\includegraphics[keepaspectratio]{/Users/andreablasco/Documents_local/Projects/01_education/03_OpenSyllabus/gender-teaching-prefs/output/figures//sim-by-cntry-1.pdf}}
\DIFaddendFL \caption{\label{fig:montecarlo-by-cntry}This figure disaggregates the comparison of proportions between actual (green) and simulated (orange) mixed-gender courses (MF/FM) by country. Simulated mixed-gender teams consistently exceed observed proportions in all countries. EU countries are Austria, Denmark, France, Germany, Ireland, Italy, Netherlands, Poland, Portugal, Spain and Sweden.}

\end{figure}

Figure \ref{fig:montecarlo-by-cntry} illustrates that the proportion of mixed-gender teams in the simulations consistently exceeds the observed proportions across various geographic regions, despite notable variations in the magnitude of these differences. In 2019, for example, mixed-gender teams are relatively rare in Great Britain (1\%) compared to Canada (4\%). However, in the simulations, the proportion of mixed-gender teams in both countries increases substantially, 1\% to 3\% and 4\% to 5\%, respectively. In 2019, the most considerable discrepancy between simulated and observed proportions is in EU countries, where it ranges from 5\% to 12\%. Overall, these patterns underscore the robustness of the findings across diverse geographic contexts.

\subsection{\DIFdelbegin \DIFdel{The Association Between Team Configuration }\DIFdelend \DIFaddbegin \DIFadd{Analysis Pre- }\DIFaddend and \DIFdelbegin \DIFdel{Course Content}\DIFdelend \DIFaddbegin \DIFadd{Post- COVID-19 in Europe}\DIFaddend }\DIFdelbegin %DIFDELCMD < \label{the-association-between-team-configuration-and-course-content}
%DIFDELCMD < %%%
\DIFdelend \DIFaddbegin \label{analysis-pre--and-post--covid-19-in-europe}
\DIFaddend 

\DIFaddbegin \DIFadd{Using a recent dataset of 2.2 million European syllabi from 2017--2022, we examine how the composition of teaching teams changed before and after the COVID-19 pandemic. Across countries and academic fields, the total number of courses declined sharply in the post-COVID period (Figure \ref{fig:covid-trend}, in SI). However, regression analyses reveal that the magnitude of this decline differed substantially depending on the gender composition of instructors (Figure \ref{fig:covid-coeffs}). After COVID-19, the share of single-instructor courses fell by more than 26\%, a significant and substantial drop relative to the pre-pandemic period. Within multi-instructor teams, the proportion of mixed-gender collaborations decreased by up to 10\%. Finally, among single-instructor courses, the share of women instructors declined by approximately 1\% relative to men. Taken together, these results highlight that single-instructor courses declined more sharply than collaborative ones, and that among collaborative teams, mixed-gender groups were disproportionately affected by the pandemic. Women were also slightly more likely to reduce their participation, although this effect is modest in size.
}
\begin{figure}
\centering
\pandocbounded{\includegraphics[keepaspectratio]{docs/report_files/figure-latex/covid-coeffs-1.pdf}}
\caption{\label{fig:covid-coeffs}\textbf{\DIFaddFL{Team composition post-Covid-19 change}}\DIFaddFL{. Coefficients and 95\% confidence intervals for post-covid period obtained from the following Binomial regressions: (a) probability of single teams being taught by a woman vs man; (b) probability of teams being composed by one instructor vs multiple (teaching collaboration); (c) probability of teaching teams being mixed-gender vs.~same-gender collaborations. All regressions are based on 2.2M syllabi from institutions in thirty-three European countries (excluding teams with three or more collaborations, and authors of unknown gender).}}
\end{figure}

\DIFadd{Figure \ref{fig:covid-simulations} shows the results of our Montecarlo simulations for this dataset. As highlighted in the previous analysis, there are disproportionately fewer mixed-gender teams than one would expect under gender-neutral matching. However, post-Covid, the gap appears to widen substantially in 2022. This pattern indicates that the pandemic, particularly the lockdown measures implemented across Europe, may have significantly hindered the formation of mixed-gender teaching teams.
}
\begin{figure}
\centering
\pandocbounded{\includegraphics[keepaspectratio]{docs/report_files/figure-latex/covid-simulations-1.pdf}}
\caption{\label{fig:covid-simulations}\textbf{\DIFaddFL{Post-Pandemic Team Formation}}\DIFaddFL{. Simulated gender-neutral vs observed teaching teams for 2.2M syllabi from institutions in thirty-three European countries. Covid-19 period highlighted (2020-2021). This shows that the gap between mixed-gender and same-gender collaborations widens after Covid.}}
\end{figure}

\subsection{\DIFadd{Associations Between Teams and Course Content}}\label{associations-between-teams-and-course-content}

\DIFaddend To examine how teaching team size and gender composition are associated with important aspects of course materials, we analyse three key metrics: the percentile rank of the interdisciplinarity score per year (``interdisciplinarity''), the average publication year of the readings (``age of readings''), and the share of female authors cited in the assigned readings (``share of female authors''). As previously discussed, we employ linear mixed-effects regressions to examine these relationships while accounting for unobserved differences across academic fields, institutions, time, and other relevant covariates.

\subsubsection{Interdisciplinarity}\label{interdisciplinarity-1}

\begin{figure}
\DIFdelbeginFL %DIFDELCMD < 

%DIFDELCMD < {%%%
\DIFdelendFL \centering
\DIFdelbeginFL %DIFDELCMD < \includegraphics[width=1\linewidth]{/Users/andreablasco/Documents_local/Work/02_Projects/01_Education_Projects/03_OpenSyllabus/gender-teaching-prefs/output/figures//plot-interdisciplinarity-1} 
%DIFDELCMD < 

%DIFDELCMD < }
%DIFDELCMD < 

%DIFDELCMD < %%%
\DIFdelendFL \DIFaddbeginFL \pandocbounded{\includegraphics[keepaspectratio]{docs/report_files/figure-latex/interdisciplinarity-1.pdf}}
\DIFaddendFL \caption{\DIFdelbeginFL \DIFdelFL{This figure illustrates the yearly difference }\DIFdelendFL \DIFaddbeginFL \label{fig:interdisciplinarity}\textbf{\DIFaddFL{Interdisciplinarity}}\DIFaddFL{. Estimates and 95\% confidence intervals for differences }\DIFaddendFL in \DIFdelbeginFL \DIFdelFL{interdisciplinarity relative to }\DIFdelendFL \DIFaddbeginFL \DIFaddFL{Interdisciplinarity between }\DIFaddendFL courses taught by \DIFdelbeginFL \DIFdelFL{a single male instructor }\DIFdelendFL \DIFaddbeginFL \DIFaddFL{men alone }\DIFaddendFL (\DIFdelbeginFL \DIFdelFL{baseline = 0}\DIFdelendFL \DIFaddbeginFL \DIFaddFL{reference}\DIFaddendFL ) \DIFdelbeginFL \DIFdelFL{for each team configuration. The values are regression coefficients obtained separately for each academic year, with controls for country }\DIFdelendFL and \DIFdelbeginFL \DIFdelFL{field}\DIFdelendFL \DIFaddbeginFL \DIFaddFL{courses taught by other instructor teams}\DIFaddendFL . Interdisciplinarity is the percentile rank\DIFaddbeginFL \DIFaddFL{, within a given year, }\DIFaddendFL of \DIFdelbeginFL \DIFdelFL{each course’s }\DIFdelendFL \DIFaddbeginFL \DIFaddFL{a syllabus' }\DIFaddendFL interdisciplinarity score \DIFdelbeginFL \DIFdelFL{within }\DIFdelendFL \DIFaddbeginFL \DIFaddFL{based on }\DIFaddendFL its \DIFdelbeginFL \DIFdelFL{cohort; positive values}\DIFdelendFL \DIFaddbeginFL \DIFaddFL{textual description relative to other syllabi. Estimates are from separate mixed-effects regressions for each academic year}\DIFaddendFL , \DIFdelbeginFL \DIFdelFL{therefore}\DIFdelendFL \DIFaddbeginFL \DIFaddFL{controlling for country}\DIFaddendFL , \DIFdelbeginFL \DIFdelFL{indicate }\DIFdelendFL \DIFaddbeginFL \DIFaddFL{field, institution, course level, probability that the document is }\DIFaddendFL a \DIFdelbeginFL \DIFdelFL{greater interdisciplinarity}\DIFdelendFL \DIFaddbeginFL \DIFaddFL{syllabus, and number of readings, with weights proportional to the square root of the number of readings}\DIFaddendFL . \DIFdelbeginFL \DIFdelFL{The shaded area represents 95\% confidence intervals}\DIFdelendFL \DIFaddbeginFL \DIFaddFL{Smoothed curves show local estimates averaged across years, weighted by the precision (standard errors) of the estimates}\DIFaddendFL .}
\DIFdelbeginFL %DIFDELCMD < \label{fig:interdisciplinarity}
%DIFDELCMD < %%%
\DIFdelendFL \end{figure}

Figure \ref{fig:interdisciplinarity} illustrates the evolving variation in interdisciplinarity across different team configurations, accounting for unobserved effects for academic field and country. When teaching alone, our analysis reveals little differences in interdisciplinarity between female and male instructors. However, in recent years, female instructors have increasingly taught significantly more interdisciplinary solo courses. By contrast, mixed-gender consistently exhibit higher interdisciplinarity compared to all-male teams. As the figure illustrates, in most years, the difference between mixed-gender and male-alone courses is statistically significant (p \textless{} 0.05), and teams with two male instructors tend to be less interdisciplinary than solo-male courses. Notably, we find no significant evidence of a difference between mixed-gender and teams with two female instructors. These results point to a robust association between gender diversity and interdisciplinarity.

\subsubsection{Citing Women Authors}\label{citing-women-authors}

\begin{figure}
\DIFdelbeginFL %DIFDELCMD < 

%DIFDELCMD < {%%%
\DIFdelendFL \centering
\DIFdelbeginFL %DIFDELCMD < \includegraphics[width=1\linewidth]{/Users/andreablasco/Documents_local/Work/02_Projects/01_Education_Projects/03_OpenSyllabus/gender-teaching-prefs/output/figures//plot-female-ratio-1} 
%DIFDELCMD < 

%DIFDELCMD < }
%DIFDELCMD < 

%DIFDELCMD < %%%
\DIFdelendFL \DIFaddbeginFL \pandocbounded{\includegraphics[keepaspectratio]{docs/report_files/figure-latex/female-authors-1.pdf}}
\DIFaddendFL \caption{\DIFdelbeginFL \DIFdelFL{This figure plots the yearly difference }\DIFdelendFL \DIFaddbeginFL \label{fig:female-authors}\textbf{\DIFaddFL{Female-Authored Radings Ratio}}\DIFaddFL{. Estimates and 95\% confidence intervals for differences }\DIFaddendFL in \DIFdelbeginFL \DIFdelFL{the share of female-authored readings relative to }\DIFdelendFL \DIFaddbeginFL \DIFaddFL{Female-authored radings ratio between }\DIFaddendFL courses taught by \DIFdelbeginFL \DIFdelFL{a single male instructor }\DIFdelendFL \DIFaddbeginFL \DIFaddFL{men alone }\DIFaddendFL (\DIFdelbeginFL \DIFdelFL{baseline = 0}\DIFdelendFL \DIFaddbeginFL \DIFaddFL{reference}\DIFaddendFL ) \DIFdelbeginFL \DIFdelFL{, for each team configuration. The values are regression coefficients obtained separately for each academic year, with controls for country }\DIFdelendFL and \DIFdelbeginFL \DIFdelFL{field}\DIFdelendFL \DIFaddbeginFL \DIFaddFL{courses taught by other instructor teams}\DIFaddendFL . \DIFdelbeginFL \DIFdelFL{The outcome }\DIFdelendFL \DIFaddbeginFL \DIFaddFL{This variable }\DIFaddendFL is the \DIFdelbeginFL \DIFdelFL{percentage-point gap in the }\DIFdelendFL proportion of female authors cited in the course readings\DIFdelbeginFL \DIFdelFL{; positive values}\DIFdelendFL \DIFaddbeginFL \DIFaddFL{. Estimates are from separate mixed-effects regressions for each academic year}\DIFaddendFL , \DIFdelbeginFL \DIFdelFL{therefore}\DIFdelendFL \DIFaddbeginFL \DIFaddFL{controlling for country}\DIFaddendFL , \DIFdelbeginFL \DIFdelFL{indicate }\DIFdelendFL \DIFaddbeginFL \DIFaddFL{field, institution, course level, probability that the document is }\DIFaddendFL a \DIFdelbeginFL \DIFdelFL{larger share }\DIFdelendFL \DIFaddbeginFL \DIFaddFL{syllabus, and number }\DIFaddendFL of \DIFdelbeginFL \DIFdelFL{female authors over total authors associated }\DIFdelendFL \DIFaddbeginFL \DIFaddFL{readings, }\DIFaddendFL with \DIFdelbeginFL \DIFdelFL{a given team configuration in that year}\DIFdelendFL \DIFaddbeginFL \DIFaddFL{weights proportional to the square root of the number of readings}\DIFaddendFL . \DIFdelbeginFL \DIFdelFL{The shaded area represents 95\% confidence intervals}\DIFdelendFL \DIFaddbeginFL \DIFaddFL{Smoothed curves show local estimates averaged across years, weighted by the precision (standard errors) of the estimates}\DIFaddendFL .}
\DIFdelbeginFL %DIFDELCMD < \label{fig:female-authors}
%DIFDELCMD < %%%
\DIFdelendFL 
\end{figure}

Figure \ref{fig:female-authors} highlights significant and persistent gender association with the percentage of female authors over total authors cited in course readings per year, controlling for academic field and country.\footnote{This analysis excludes syllabi with no matched readings with the available bibliographic sources and, within the matched readings, excludes the references where the authors' gender remained unidentified.} All-female courses consistently cite a higher fraction of female authors than all-male ones in the same year, with an effect that decreases over time, going from \DIFdelbegin \DIFdel{6}\DIFdelend \DIFaddbegin \DIFadd{8}\DIFaddend \% in 2000 to \DIFdelbegin \DIFdel{3}\DIFdelend \DIFaddbegin \DIFadd{4}\DIFaddend \% in 2019, with minor or insignificant differences between courses taught by one vs.~two instructors. Mixed-gender courses are in between, as they cite a larger share of women than all-male courses but less than all-female courses. Overall, these results indicate a consistent association between gender and cited patterns, suggesting that readings tend to cluster based on the gender of the instructors, and mixed-gender teams could promote greater diversity in assigned readings.

\subsubsection{\DIFaddbegin \DIFadd{Teaching Conventionality, }\DIFaddend Age of Readings\DIFdelbegin \DIFdel{(``novelty'')}\DIFdelend \DIFaddbegin \DIFadd{, and Scientific Novelty}\DIFaddend }\DIFdelbegin %DIFDELCMD < \label{age-of-readings-novelty}
%DIFDELCMD < %%%
\DIFdelend \DIFaddbegin \label{teaching-conventionality-age-of-readings-and-scientific-novelty}
\DIFaddend 

\DIFdelbegin %DIFDELCMD < \begin{figure}
%DIFDELCMD < 

%DIFDELCMD < {\centering \includegraphics[width=1\linewidth]{/Users/andreablasco/Documents_local/Work/02_Projects/01_Education_Projects/03_OpenSyllabus/gender-teaching-prefs/output/figures//plot-age-of-references-1} 
%DIFDELCMD < 

%DIFDELCMD < }
%DIFDELCMD < 

%DIFDELCMD < %%%
%DIFDELCMD < \caption{%
{%DIFAUXCMD
\DIFdelFL{This figure illustrates the yearly difference in the percentile rank of the Age of References relative to courses taught by a single male instructor (baseline = 0) for each team configuration. The Age of Reference variable is the syllabus’ year minus the average publication year of the assigned readings. The coefficients are estimated separately for each academic year, controlling for country and broad field. Positive values indicate that the team configuration assigns newer readings compared to the male-alone baseline in the same year. The shaded area represents 95\% confidence intervals.}}%DIFAUXCMD
%DIFDELCMD < \label{fig:age-readings}
%DIFDELCMD < \end{figure}
%DIFDELCMD < 

%DIFDELCMD < %%%
\DIFdelend Figure \ref{fig:age-readings} reveals \DIFaddbegin \DIFadd{persistent temporal trends and }\DIFaddend significant gender differences in the publication age of the selected readings\DIFaddbegin \DIFadd{, even }\DIFaddend after accounting for \DIFdelbegin \DIFdel{differences across fields }\DIFdelend \DIFaddbegin \DIFadd{observable differences across multiple dimensions, including course level, academic field, institution, }\DIFaddend and country. \DIFdelbegin \footnote{\DIFdel{This analysis includes only courses with matched bibliographic data.}} %DIFAUXCMD
\addtocounter{footnote}{-1}%DIFAUXCMD
\DIFdel{Between 2000 and 2019, courses }\DIFdelend \DIFaddbegin \DIFadd{Course }\DIFaddend taught by a single female instructor tend to assign \DIFdelbegin \DIFdel{newer readings than otherwise similar }\DIFdelend \DIFaddbegin \DIFadd{more recent readings than comparable }\DIFaddend courses taught by a single man \DIFdelbegin \DIFdel{in that year, with a difference that goes }\DIFdelend \DIFaddbegin \DIFadd{instructor (the reference category), with the difference going }\DIFaddend from -3\% in 2000 to -1\% in 2019. Similarly, courses led by two women also assign newer material than those taught by two men\DIFdelbegin \DIFdel{, with a decreasing difference}\DIFdelend . Mixed-gender teams \DIFdelbegin \DIFdel{tend to assign newer readings than the baseline in a few years but not consistently across all years. There are no significant differences }\DIFdelend \DIFaddbegin \DIFadd{generally assign more recent readings than single-male instructors, although less than all-female teams of equal size. No significant differences emerge }\DIFaddend between courses taught by two male instructors and the single-male baseline. Overall, these findings \DIFdelbegin \DIFdel{suggest that women tend to assign more novel }\DIFdelend \DIFaddbegin \DIFadd{indicate that the presence of women in teaching teams, either alone or with a male co-instructor, is associated with more recent }\DIFaddend readings, although \DIFdelbegin \DIFdel{this gender difference tends }\DIFdelend \DIFaddbegin \DIFadd{the gender effect appears }\DIFaddend to decrease over time.

\DIFdelbegin \subsection{\DIFdel{Robustness}}%DIFAUXCMD
\addtocounter{subsection}{-1}%DIFAUXCMD
%DIFDELCMD < \label{robustness}
%DIFDELCMD < %%%
\DIFdelend \DIFaddbegin \begin{figure}
\centering
\pandocbounded{\includegraphics[keepaspectratio]{docs/report_files/figure-latex/age-readings-1.pdf}}
\caption{\label{fig:age-readings}\textbf{\DIFaddFL{Age of Readings by Team}}\DIFaddFL{. Estimates and 95\% confidence intervals for differences in Age of Readings between courses taught by men alone (reference) and courses taught by other instructor teams. Age of Readings is the percentile rank, within a given year, of a syllabus' average publication age of the readings. Estimates are from separate mixed-effects regressions for each academic year, controlling for country, field, institution, course level, probability that the document is a syllabus, and number of readings, with weights proportional to the square root of the number of readings. Smoothed curves show local estimates averaged across years, weighted by the precision (standard errors) of the estimates.}}
\end{figure}
\DIFaddend 

\DIFdelbegin \DIFdel{To ensure the reliability and robustness of our results, we conducted several additional analyses. In addition to reporting results from separate regressions on subsets of data by year, we also estimated mixed-effects regressions using the full dataset. These results are broadly consistent with those presented in the paper.
}\DIFdelend \DIFaddbegin \DIFadd{Figure \ref{fig:conventionality} illustrates significant gender differences in the readings' conventionality, which is the median citations of syllabus readings from references in other syllabi. Team-taught courses tend to be more conventional than single-instructor courses, as both mixed-gender and all-male teams exhibit significantly higher conventionality than the baseline, while the difference with two-women teams is insignificant. Courses taught by women alone are significantly less conventional than those taught by men alone, but the difference is small (less than 1 pp). Taken together, this evidence indicates a strong association between reading conventionality and teaching in teams with more than one instructor, which is partly offset by the presence of women.
}\DIFaddend 

\DIFdelbegin \DIFdel{Furthermore, while the reported model specification is our preferred one , we explored several alternatives.
First, we log-transformed our dependent variables -- the age of readings, interdisciplinarity score,and count of female authors --- to explore potential non-linear relationships, particularly relevant for disciplines with old readings (e. g., history)
}\DIFdelend \DIFaddbegin \begin{figure}
\centering
\pandocbounded{\includegraphics[keepaspectratio]{docs/report_files/figure-latex/conventionality-1.pdf}}
\caption{\label{fig:conventionality}\textbf{\DIFaddFL{Readings Conventionality by Team}}\DIFaddFL{. Estimates and 95\% confidence intervals for differences in Teaching Conventionality between courses taught by men alone (reference) and courses taught by other instructor teams. Teaching Conventionality is the percentile rank, within a given year, of a syllabus' median citations of readings, based on references from other syllabi. Estimates are from separate mixed-effects regressions for each academic year, controlling for country, field, institution, course level, probability that the document is a syllabus, and number of readings, with weights proportional to the square root of the number of readings. Smoothed curves show local estimates averaged across years, weighted by the precision (standard errors) of the estimates.}}

\end{figure}
\DIFaddend 

\DIFdelbegin \DIFdel{Second, we employed a quasi-Poisson mixed-effect model instead of the reported linear mixed-effect regression to better account for non-linearities in the count of female authors.We also trimmed the data to exclude outliers, such as courseswith an unusually high number of readings. These changes yield similar results}\DIFdelend \DIFaddbegin \DIFadd{Figure \ref{fig:atypicality} presents the results of our analysis of the scientific novelty of the readings, based on approximately 400,000 syllabi for which we obtained Uzzi et al. (\mbox{%DIFAUXCMD
\citeproc{ref-uzzi2013atypical}{47}}\hskip0pt%DIFAUXCMD
)'s atypicality measure precomputed by SciSciNet (v2) (\mbox{%DIFAUXCMD
\citeproc{ref-lin2023sciscinet}{46}}\hskip0pt%DIFAUXCMD
). Controlling for field, course level, and other relevant characteristics, the results indicate that readings produced by mixed-gender teams are significantly less at the frontier of knowledge.}\footnote{
  \DIFadd{This regression also inclused the age of readings as a control for the novelty at time of publication of the article.}} \DIFadd{Similarly, women-led teams show a negative association compared to man-teaching-alone courses, with the effect being stronger in courses taught by two women. In other words, men-only teams exhibit slightly higher atypicality in our dataset than any other team configuration}\DIFaddend .

\DIFdelbegin \DIFdel{We further tested alternative model specifications, such as considering the number of instructors as a numeric variable interacted with the proportion of female instructors to explore a linear association with the dependent variables. While we observe similar results, the linear association could be misleading when extrapolated beyond two instructors. So, we preferred to keep team configurations as fixed effects, as reported here. }\DIFdelend \DIFaddbegin \begin{figure}
\centering
\pandocbounded{\includegraphics[keepaspectratio]{docs/report_files/figure-latex/atypicality-1.pdf}}
\caption{\label{fig:atypicality}\textbf{\DIFaddFL{Scientific Novelty by Team}}\DIFaddFL{. Estimates and 95\% confidence intervals for differences in Scientific Novelty between courses taught by men alone (reference) and courses taught by other instructor teams. Scientific Novelty is the percentile rank, within a given year, of a syllabus' median Uzzi et al's atypicality of readings, based on the combination of references in the readings. Estimates are from separate mixed-effects regressions for each academic year, controlling for country, field, institution, course level, probability that the document is a syllabus, and number of readings, with weights proportional to the square root of the number of readings. Smoothed curves show local estimates averaged across years, weighted by the precision (standard errors) of the estimates.}}
\end{figure}
\DIFaddend 

\DIFdelbegin \DIFdel{Finally, we further explored the specifications that distinguish between male-led and female-led courses, using the gender of the first listed instructor as a proxy for seniority}\DIFdelend \DIFaddbegin \subsection{\DIFadd{Heterogeneity}}\label{heterogeneity}

\DIFadd{Figure \ref{fig:stem} highlights several key differences between STEM and non-STEM fields, based on separate models estimated for STEM fields compared against results from the full sample.}\footnote{
  \DIFadd{We converted CIP fields into ISCED-F broad categories and considered STEM the following: (1) Natural Sciences, Mathematics, and Statistics,
  (2) Information and Communication Technologies (ICTs), and (3) Engineering, Manufacturing, and Construction.}}\DIFadd{. The gap in female-author citations associated with women instructors remains consistent but is significantly smaller in STEM fields, roughly half the size observed in all fields. The positive association between mixed-gender teams and conventional readings is larger in STEM fields. Conversely, the gap in atypicality widens for women teaching alone, although the effect has been diminishing over time. In contrast, the gap in the age of assigned readings between women and men teaching alone is smaller in STEM. With respect to interdisciplinarity, the estimated differences are small and not consistently significant across years. As a result, the evidence for heterogeneity in interdisciplinarity within STEM is inconclusive.
}
\begin{figure}
\centering
\pandocbounded{\includegraphics[keepaspectratio]{docs/report_files/figure-latex/stem-1.pdf}}
\caption{\label{fig:stem}\textbf{\DIFaddFL{Heterogeneity within STEM Fields}}\DIFaddFL{. Estimates and 95\% confidence intervals for differences in (1) interdisciplinarity (2) readings' conventionality (3) age of readings (4) female-authored ratio and (5) readings' atypicality. Estimates are from separate mixed-effects regressions for STEM fields vs the full sample, in each academic year, controlling for country, field, institution, course level, probability that the document is a syllabus, and number of readings, with weights proportional to the square root of the number of readings. Smoothed curves show local estimates averaged across years, weighted by the precision (standard errors) of the estimates.}}

\end{figure}

\DIFadd{We have also explored heterogeneity across countries by fitting separate models for the US and comparing the results against the estimates from the full sample. The analysis revealed no signficant differences for the US-based regressions (Figure \ref{fig:hetero-us-plots}, SI). Therefore, our study provides no evidence of cross-country differences in the observed associations.
}

\subsection{\DIFadd{Robustness}}\label{robustness}

\DIFadd{The robustness checks conducted, as detailed in the Supplementary Information (Section \ref{robustness}), show that the main results are consistent across multiple alternative analyses. These included mixed-effects regressions on the full dataset, log-transformations of key dependent variables, quasi-Poisson models for counts of female authors, trimming outliers, and testing different specifications related to instructor composition and seniority. Across all these variations, the findings remained broadly stable, demonstrating that the results are not sensitive to specific modeling choices or data subsets}\DIFaddend .
\DIFdelbegin \DIFdel{Overall, results appear robust to these additional analyses.
}\DIFdelend 

\section{Discussion}\label{discussion}

\subsection{Summary of the findings}\label{summary-of-the-findings}

Our analysis of approximately five million syllabi from over 4,000 universities in fifteen countries reveals several key findings. First, while the share of classes taught by female instructors has been increasing over the last twenty years, the fraction of mixed-gender co-teaching classes has remained consistently low (3\%) and significantly below that of same-gender teams (10\%). This trend can stem from a lack of gender diversity in specific fields or \DIFaddbegin \DIFadd{from }\DIFaddend administrative practices that limit opportunities \DIFdelbegin \DIFdel{for forming }\DIFdelend \DIFaddbegin \DIFadd{to form }\DIFaddend mixed-gender collaborations. However, it can also arise from gender preferences in team formation, whereby instructors tend to form teams with same-gender partners (``homophily''), a tendency already observed in research collaborations and other academic domains (e.g., \DIFdelbegin \DIFdel{\mbox{%DIFAUXCMD
\citeproc{ref-kwiek2021gender}{17}}\hskip0pt%DIFAUXCMD
}\DIFdelend \DIFaddbegin \DIFadd{\mbox{%DIFAUXCMD
\citeproc{ref-kwiek2021gender}{18}}\hskip0pt%DIFAUXCMD
}\DIFaddend ).

\DIFdelbegin \DIFdel{Moreover, our }\DIFdelend \DIFaddbegin \DIFadd{Our }\DIFaddend analysis shows that the underrepresentation of mixed-gender teams persists even after accounting for unobserved differences across fields, institutions, and \DIFaddbegin \DIFadd{academic }\DIFaddend years. Specifically, employing Monte Carlo simulations that randomly form teams in a gender-neutral manner, our analysis controls for typical institutional constraints, such as maintaining constant course loads across institutions and 69 academic fields within a given year. By comparing simulated to actual teams, we find that instructors tend to partner with same-gender colleagues more than twice as expected between 2005 and 2019. This result is robust across academic fields, although it is more pronounced in specific fields (\DIFdelbegin \DIFdel{e.g., }\DIFdelend Chemistry, Business, Medicine). Additionally, the underrepresentation of mixed-gender teams is consistent across \DIFdelbegin \DIFdel{different }\DIFdelend countries, with the \DIFdelbegin \DIFdel{most substantial }\DIFdelend \DIFaddbegin \DIFadd{greatest }\DIFaddend effect in EU countries, where only 5\% of courses are taught by mixed-gender teams, compared to an expected 12\%.

\DIFaddbegin \DIFadd{Additional insights emerged from the analysis of team patterns before and after the COVID-19 pandemic (2017-2022) across thirty-three European countries. First, the pandemic caused a general decline in the number of courses, consistent across fields and countries, but this decline was disproportionately concentrated among individual instructors rather than collaborative teams. This finding likely reflects the reduction in in-person course offerings, which were primarily taught individually. Second, although collaborative teams were relatively less affected by the overall drop in courses, mixed-gender teams declined disproportionately. Taken together, these findings suggest that the formation of mixed-gender teaching teams depends more on in-person interactions than other types of teaching collaborations.
}

\DIFadd{Our analysis further highlights how different team configurations are associated with key aspects of teaching: (i) mixed-gender teams tend to be more interdisciplinary than all-male teams, (ii) all-female teams tend to assign more recent readings, (iii) mixed-gender teams tend to cite more ``conventional'' readings (i.e., those cited in other syllabi before) and less at the frontier of scientific research (i.e., lower mean atypicality) and (iv) mixed-gender teams tend to cite a higher share of female authors than all-male teams, but less than all-female teams. These associations are not driven by unobserved differences across fields, course levels, institutions, and years that we account for in the regressions. Furthermore, these findings expand existing literature in the following ways.
}

\DIFaddend \subsection{Main Contributions}\label{main-contributions}

\DIFdelbegin \DIFdel{These findings }\DIFdelend \DIFaddbegin \DIFadd{First, our findings on team formation }\DIFaddend extend prior research on academic collaboration, which has traditionally focused on research partnerships (\DIFdelbegin \DIFdel{\mbox{%DIFAUXCMD
\citeproc{ref-guimera2005team}{11}}\hskip0pt%DIFAUXCMD
, \mbox{%DIFAUXCMD
\citeproc{ref-wuchty2007increasing}{12}}\hskip0pt%DIFAUXCMD
, \mbox{%DIFAUXCMD
\citeproc{ref-yang2022gender}{14}}\hskip0pt%DIFAUXCMD
, \mbox{%DIFAUXCMD
\citeproc{ref-boschini2007team}{15}}\hskip0pt%DIFAUXCMD
, \mbox{%DIFAUXCMD
\citeproc{ref-maddi2021gender}{31}}\hskip0pt%DIFAUXCMD
}\DIFdelend \DIFaddbegin \DIFadd{\mbox{%DIFAUXCMD
\citeproc{ref-wuchty2007increasing}{12}}\hskip0pt%DIFAUXCMD
, \mbox{%DIFAUXCMD
\citeproc{ref-boschini2007team}{17}}\hskip0pt%DIFAUXCMD
, \mbox{%DIFAUXCMD
\citeproc{ref-yang2022gender}{32}}\hskip0pt%DIFAUXCMD
, \mbox{%DIFAUXCMD
\citeproc{ref-guimera2005team}{53}}\hskip0pt%DIFAUXCMD
, \mbox{%DIFAUXCMD
\citeproc{ref-maddi2021gender}{54}}\hskip0pt%DIFAUXCMD
}\DIFaddend ), to the domain of teaching. The significant gender patterns and underrepresentation of mixed-gender co-teaching teams echo the literature on homophily and bias in team formation within research collaborations (e.g., \DIFdelbegin \DIFdel{\mbox{%DIFAUXCMD
\citeproc{ref-kwiek2021gender}{17}}\hskip0pt%DIFAUXCMD
}\DIFdelend \DIFaddbegin \DIFadd{\mbox{%DIFAUXCMD
\citeproc{ref-kwiek2021gender}{18}}\hskip0pt%DIFAUXCMD
}\DIFaddend ) and underscore the critical role of gender dynamics \DIFdelbegin \DIFdel{not only in }\DIFdelend \DIFaddbegin \DIFadd{in shaping not only }\DIFaddend scientific production but also \DIFdelbegin \DIFdel{in shaping }\DIFdelend pedagogical practices. Moreover, by combining \DIFdelbegin \DIFdel{the }\DIFdelend analysis of a massive syllabi dataset with Monte Carlo simulations of \DIFdelbegin \DIFdel{a }\DIFdelend gender-neutral team formation, our findings show that these patterns cannot be explained solely by country-level, institutional\DIFaddbegin \DIFadd{, }\DIFaddend and field-specific factors, such as cultural barriers or \DIFdelbegin \DIFdel{a }\DIFdelend gender imbalance in specific fields. This evidence highlights the need to investigate these gender dynamics further, focusing on the underlying causal mechanisms that limit gender-diverse collaborations in teaching, a critical direction for future research.

\DIFdelbegin \DIFdel{Our analysis further highlights how different team configurations are associated with key aspects of teaching: (1) mixed-gender teams tend to be more interdisciplinary than all-male teams, (2) all-female teams tend to assign more novel readings, and (3) mixed-gender teams tend to cite a higher share of female authors than all-male teams, but less than all-female teams. These associations are not driven by unobserved differences across fields, institutions, and years that we account for in the regressions. Furthermore, these findings expand existing literature in the following ways.
}%DIFDELCMD < 

%DIFDELCMD < %%%
\DIFdel{First, }\DIFdelend \DIFaddbegin \DIFadd{Secondly, }\DIFaddend building on prior research that measures interdisciplinary curricula (\DIFdelbegin \DIFdel{\mbox{%DIFAUXCMD
\citeproc{ref-evans2016measuring}{26}}\hskip0pt%DIFAUXCMD
}\DIFdelend \DIFaddbegin \DIFadd{\mbox{%DIFAUXCMD
\citeproc{ref-evans2016measuring}{51}}\hskip0pt%DIFAUXCMD
}\DIFaddend ) and examines their effects on student outcomes (\citeproc{ref-han2023interdisciplinary}{3}), our study shifts the focus to the drivers of interdisciplinarity, linking it to the size and gender composition of teaching teams. While we find only minor and inconsistent gender differences in solo-taught teams, interdisciplinarity significantly increases in mixed-gender teaching teams. This effect is associated with the presence of women collaborators, who tend to be more interdisciplinary when working in teams. These findings extend the literature that compares interdisciplinarity between men and women in academic research (\DIFdelbegin \DIFdel{\mbox{%DIFAUXCMD
\citeproc{ref-pinheiro2022women}{39}}\hskip0pt%DIFAUXCMD
}\DIFdelend \DIFaddbegin \DIFadd{\mbox{%DIFAUXCMD
\citeproc{ref-pinheiro2022women}{30}}\hskip0pt%DIFAUXCMD
}\DIFaddend ) to the teaching domain, establishing a novel connection between gender diversity and interdisciplinary teaching practices. This connection highlights a potential lever for universities and academic institutions seeking to foster interdisciplinarity in their curricula.

\DIFdelbegin \DIFdel{Secondly}\DIFdelend \DIFaddbegin \DIFadd{Thirdly}\DIFaddend , our study advances our understanding of the determinants of gender disparities in academic recognition \DIFdelbegin %DIFDELCMD < {[}%%%
\DIFdelend by extending the focus beyond scholarly citations to teaching materials. While prior research has \DIFdelbegin \DIFdel{extensively }\DIFdelend \DIFaddbegin \DIFadd{predominantly }\DIFaddend focused on gender citation gaps in academic publications \DIFdelbegin %DIFDELCMD < {[}%%%
\DIFdel{REFs}%DIFDELCMD < {]}%%%
\DIFdel{, our findings illustrate }\DIFdelend \DIFaddbegin \DIFadd{(e.g., \mbox{%DIFAUXCMD
\citeproc{ref-dion2018gendered}{21}}\hskip0pt%DIFAUXCMD
), our study demonstrates }\DIFaddend the critical role of gender dynamics in the content presented in class and assigned readings. \DIFdelbegin \DIFdel{These findings }\DIFdelend \DIFaddbegin \DIFadd{More specifically, our findings align with prior research on gendered citation patterns in the academic literature (e.g., \mbox{%DIFAUXCMD
\citeproc{ref-dion2018gendered}{21}}\hskip0pt%DIFAUXCMD
), which shows that female scholars are more likely to cite work by other women. Interestingly, Dion et al. (\mbox{%DIFAUXCMD
\citeproc{ref-dion2018gendered}{21}}\hskip0pt%DIFAUXCMD
) finds that mixed-gender research teams do not cite more women than male-only teams, highlighting a divergence from our course-level results. This finding indicates that instructor gender shapes reading lists in ways that may differ from research citation patterns and that mixed-gender teaching teams could foster greater gender diversity in assigned readings. Together, these findings }\DIFaddend suggest new potential drivers of gender stereotypes in academia (\DIFdelbegin \DIFdel{\mbox{%DIFAUXCMD
\citeproc{ref-chatterjee2021gender}{34}}\hskip0pt%DIFAUXCMD
}\DIFdelend \DIFaddbegin \DIFadd{\mbox{%DIFAUXCMD
\citeproc{ref-chatterjee2021gender}{49}}\hskip0pt%DIFAUXCMD
}\DIFaddend ), underscoring the need to broaden research on the causal mechanisms underlying gender disparities in citation practices to include teaching contexts, thereby moving beyond the current emphasis on scholarly citations alone.

Finally, our findings reveal a persistent gender-based \DIFdelbegin \DIFdel{association with }\DIFdelend \DIFaddbegin \DIFadd{pattern in }\DIFaddend the assignment of novel materials, \DIFdelbegin \DIFdel{whereby female instructors tend }\DIFdelend \DIFaddbegin \DIFadd{with female instructors tending }\DIFaddend to include more \DIFdelbegin \DIFdel{novel }\DIFdelend \DIFaddbegin \DIFadd{recent }\DIFaddend content in their syllabi than male instructors. \DIFdelbegin \DIFdel{This trend holds even when accounting for the order of appearance in the syllabus, which serves as a proxy for the instructional experience or hierarchical position. This evidence }\DIFdelend \DIFaddbegin \DIFadd{In contrast, mixed-gender teams---or teams in general---tend to assign more conventional readings and cite scientific articles with lower atypicality. This evidence points to a complex relationship between team composition and teaching innovation that warrants further investigation. One possible explanation is that female instructors may assign more recent readings because they enter the profession more recently, which could account for a weakening association between instructor gender and recency over time. Similarly, mixed-gender teams may rely on less atypical or more conventional readings because instructors from different fields must establish a shared foundation rooted in classic works. Taken together with our previous finding that interdisciplinary content is more common in mixed-gender teams, these results suggest that interdisciplinary courses may inherently require more conventional readings to bridge disciplinary differences.
}

\DIFadd{Overall, unveiling these associations }\DIFaddend contributes to understanding the \DIFdelbegin \DIFdel{drivers of the }\DIFdelend observed gender-innovation gap in tertiary education (e.g., \citeproc{ref-biasi2022education}{10}). \DIFdelbegin \DIFdel{While }\DIFdelend \DIFaddbegin \DIFadd{However, while }\DIFaddend our results suggest a strong gender association, they remain inconclusive regarding its underlying drivers, whether the pattern reflects gender per se or is confounded by age or career stage (e.g., younger vs.~female instructors). Notably, this trend coincides with the growing presence of women entering academia during the study period, suggesting that newer instructors may play a key role in pushing curricular boundaries. Nonetheless, further research is needed to disentangle the influence of gender from other intersecting factors, such as cohort effects, field-specific norms, and institutional context.

\subsection{Limitations}\label{limitations}

Our analysis bears several limitations. First, \DIFdelbegin \DIFdel{we lack more }\DIFdelend \DIFaddbegin \DIFadd{the dataset does not include }\DIFaddend granular information on the course's subfields that go beyond the 69 top-level fields identified in the dataset. If male and female instructors tend to concentrate on different subfields (e.g., Econometrics and Macroeconomics, within Economics), the observed underrepresentation of mixed-gender teams may reflect limited availability of instructors of the other gender in the subfield, especially for advanced courses where teaching teams are more likely to form around narrower areas of expertise. This limitation prevents disentangling the exact driver of mixed-gender underrepresentation, which can derive from gender-based preferences towards subfields, limiting availability, or a tendency to seek same-gender partners or homophily. Further research is needed to \DIFaddbegin \DIFadd{fully }\DIFaddend disentangle these underlying mechanisms\DIFdelbegin \DIFdel{fully}\DIFdelend .

\DIFdelbegin \DIFdel{Similarly, another limitation is the lack of data on potential determinants of team formation }\DIFdelend \DIFaddbegin \DIFadd{Second, our data lack characteristics }\DIFaddend that operate within institutions, such as instructor experience or academic rank. Hierarchical differences or administrative structures may contribute to the observed gap in cross-gender collaborations --- for example, \DIFdelbegin \DIFdel{if }\DIFdelend \DIFaddbegin \DIFadd{by making }\DIFaddend newly appointed female instructors \DIFdelbegin \DIFdel{are }\DIFdelend less likely to be selected by higher-ranking male instructors. However, our results \DIFdelbegin \DIFdel{illustrate }\DIFdelend \DIFaddbegin \DIFadd{show }\DIFaddend that the under-representation of mixed-gender teams persists even when co-teaching teams are disaggregated by \DIFdelbegin \DIFdel{male-led }\DIFdelend \DIFaddbegin \DIFadd{male- }\DIFaddend and female-led courses (as measured by \DIFdelbegin \DIFdel{the order of appearance in the syllabus }\DIFdelend \DIFaddbegin \DIFadd{syllabus order}\DIFaddend ), demonstrating that the trend holds regardless of hierarchical differences. Therefore, while we cannot entirely exclude other institutional constraints, the observed gap in mixed-gender teams likely reflects strong gender-based homophily.

\DIFdelbegin \DIFdel{Our findings underscore a significant and consistent underrepresentation of gender-mixed teaching teams, carrying profound implications. Limited }\DIFdelend \DIFaddbegin \DIFadd{Third, our measurement of novelty in syllabi is constrained by the limited citation base typical of course materials and the temporal mismatch between publication-time novelty and the novelty students experience during instruction. These factors limit our capacity to apply systematically standard novelty metrics to teaching materials (e.g., \mbox{%DIFAUXCMD
\citeproc{ref-uzzi2013atypical}{47}}\hskip0pt%DIFAUXCMD
). Yet, by combining multiple measures, our work provides a viable path to characterise the multidimensional nature of innovation in teaching using syllabi data.
}

\DIFadd{Fourth, institutional policies regarding teaching loads and the recognition of co-teaching vary substantially across countries and institutions. In many universities, team teaching does not fully count toward teaching loads or is treated as an overload, limiting faculty opportunities to engage in it even if they wish to. Moreover, teaching teams are not always determined by individual instructors' preferences: in many institutions, course assignments and team composition are partially or entirely determined at the departmental or program level---for example, by course coordinators, curriculum structures, or administrative rules. In such cases, team formation may reflect institutional or organizational preferences rather than instructors' choices. Our data do not capture these institutional constraints or assignment mechanisms, which may mechanically reduce or shape the observed incidence of }\DIFaddend mixed-gender \DIFdelbegin \DIFdel{collaborations in teaching may limit the exchange of information among instructors of different genders and restrict students' access to diverse courses. These dynamics, in turn, could hinder interdisciplinary engagement or reinforce structural disadvantages for women within academic networks. Our work thus identifies a potential additional driver of gender bias, complementing existing research on gender disparities in various academic settings, including research collaboration, access to funding, and student gender stereotypes}\DIFdelend \DIFaddbegin \DIFadd{teams. Consequently, our descriptive patterns reflect both individual preferences and the structural opportunities and assignment practices provided by institutions.
}

\DIFadd{Finally, although our data covers 50\% of educational institutions in the selected countries, and comparisons with official statistics on the distribution of academic staff and US graduates by field show a reasonable alignment with the trends observed in our sample, it is important to note that the sample is not fully representative of the underlying population. Additional research is needed to validate our results with appropriate weights or samples of the full population}\DIFaddend .

\subsection{Conclusion}\label{conclusion}

Overall, this study offers critical insights into the academic market. As for research teams, teaching teams can impact instructors in multiple ways. First, junior faculty may have access to the experience and advice of senior faculty, which could increase their chances of promotion. Secondly, diverse teaching teams may inspire research collaborations, including interdisciplinary work. Finally, teaching assistants can actively seek out mentors and may benefit from these collaborations when applying for PhD programs or entering the labour market. Therefore, understanding how teaching teams are formed can help explain career trajectories.

\DIFdelbegin %DIFDELCMD < \newpage
%DIFDELCMD < %%%
\DIFdelend \DIFaddbegin \clearpage
\DIFaddend \processdelayedfloats
\DIFaddbegin \clearpage
\DIFaddend 

\DIFaddbegin \section*{\DIFadd{Acknowledgements}}\label{acknowledgements}
\addcontentsline{toc}{section}{\DIFadd{Acknowledgements}}

\DIFadd{The contents of this publication do not necessarily reflect the position or opinion of the European Commission or the Joint Research Centre.
}

\section*{\DIFadd{Data Availability Statement}}\label{data-availability-statement}
\addcontentsline{toc}{section}{\DIFadd{Data Availability Statement}}

\DIFadd{All the datasets generated during the current study are available online in the }{[}\DIFadd{name of repository}{]}\DIFadd{, }{[}\DIFadd{persistent web link or DOI}{]}\DIFadd{. The raw data will be available from the corresponding author upon reasonable request and if the request meets the contractual agreements with OpenSyllabus.
}

\DIFadd{All the scripts and code used in the analyses is available at }{[}\DIFadd{GitHub link or DOI}{]}\DIFadd{. This includes two repositories one with all the Python scripts used to compute the interdisciplinarity score }{[}\DIFadd{LINK or DOI}{]} \DIFadd{and another with Rscripts to analyse the data and create the figures }{[}\DIFadd{LINK or DOI}{]}\DIFadd{.
}

\section*{\DIFadd{References}}\label{references}
\addcontentsline{toc}{section}{\DIFadd{References}}

\phantomsection\label{refs}
\begin{CSLReferences}{0}{1}
\bibitem[\citeproctext]{ref-bakkenes2010teacher}
\CSLLeftMargin{1. }%DIF > 
\CSLRightInline{I. Bakkenes, J. D. Vermunt, T. Wubbels, Teacher learning in the context of educational innovation: Learning activities and learning outcomes of experienced teachers. \emph{Learning and instruction} \textbf{20}, 533--548 (2010).}

\bibitem[\citeproctext]{ref-vembye2024effects}
\CSLLeftMargin{2. }%DIF > 
\CSLRightInline{M. H. Vembye, F. Weiss, B. Hamilton Bhat, The effects of co-teaching and related collaborative models of instruction on student achievement: A systematic review and meta-analysis. \emph{Review of Educational Research} \textbf{94}, 376--422 (2024).}

\bibitem[\citeproctext]{ref-han2023interdisciplinary}
\CSLLeftMargin{3. }%DIF > 
\CSLRightInline{S. Han, J. LaViolette, C. Borkenhagen, W. McAllister, P. S. Bearman, Interdisciplinary college curriculum and its labor market implications. \emph{Proceedings of the National Academy of Sciences} \textbf{120}, e2221915120 (2023).}

\bibitem[\citeproctext]{ref-ellemers2018gender}
\CSLLeftMargin{4. }%DIF > 
\CSLRightInline{N. Ellemers, Gender stereotypes. \emph{Annual review of psychology} \textbf{69}, 275--298 (2018).}

\bibitem[\citeproctext]{ref-lockwood2006someone}
\CSLLeftMargin{5. }%DIF > 
\CSLRightInline{P. Lockwood, {``Someone like me can be successful''}: Do college students need same-gender role models? \emph{Psychology of women quarterly} \textbf{30}, 36--46 (2006).}

\bibitem[\citeproctext]{ref-canaan2023impact}
\CSLLeftMargin{6. }%DIF > 
\CSLRightInline{S. Canaan, P. Mouganie, The impact of advisor gender on female students' STEM enrollment and persistence. \emph{Journal of Human Resources} \textbf{58}, 593--632 (2023).}

\bibitem[\citeproctext]{ref-buckingham2021learning}
\CSLLeftMargin{7. }%DIF > 
\CSLRightInline{L. R. Buckingham, A. López-Hernández, B. Strotmann, Learning by comparison: The benefits of co-teaching for university professors' professional development in \emph{Frontiers in Education}, (Frontiers Media SA, 2021), p. 776991.}

\bibitem[\citeproctext]{ref-nowogrodzki2016mining}
\CSLLeftMargin{8. }%DIF > 
\CSLRightInline{A. Nowogrodzki, Mining the secrets of college syllabuses. \emph{Nature} \textbf{539}, 125--126 (2016).}

\bibitem[\citeproctext]{ref-gao2023quantifying}
\CSLLeftMargin{9. }%DIF > 
\CSLRightInline{J. Gao, D. Wang, Quantifying the benefit of artificial intelligence for scientific research. \emph{arXiv preprint arXiv:2304.10578} (2023).}

\bibitem[\citeproctext]{ref-biasi2022education}
\CSLLeftMargin{10. }%DIF > 
\CSLRightInline{B. Biasi, S. Ma, {``The education-innovation gap''} (National Bureau of Economic Research, 2022).}

\bibitem[\citeproctext]{ref-myers2020unequal}
\CSLLeftMargin{11. }%DIF > 
\CSLRightInline{K. R. Myers, \emph{et al.}, Unequal effects of the COVID-19 pandemic on scientists. \emph{Nature human behaviour} \textbf{4}, 880--883 (2020).}

\bibitem[\citeproctext]{ref-wuchty2007increasing}
\CSLLeftMargin{12. }%DIF > 
\CSLRightInline{S. Wuchty, B. F. Jones, B. Uzzi, The increasing dominance of teams in production of knowledge. \emph{Science} \textbf{316}, 1036--1039 (2007).}

\bibitem[\citeproctext]{ref-milojevic2014principles}
\CSLLeftMargin{13. }%DIF > 
\CSLRightInline{S. Milojević, Principles of scientific research team formation and evolution. \emph{Proceedings of the National Academy of Sciences} \textbf{111}, 3984--3989 (2014).}

\bibitem[\citeproctext]{ref-mcpherson2001birds}
\CSLLeftMargin{14. }%DIF > 
\CSLRightInline{M. McPherson, L. Smith-Lovin, J. M. Cook, Birds of a feather: Homophily in social networks. \emph{Annual review of sociology} \textbf{27}, 415--444 (2001).}

\bibitem[\citeproctext]{ref-mcdowell1992effect}
\CSLLeftMargin{15. }%DIF > 
\CSLRightInline{J. M. McDowell, J. K. Smith, The effect of gender-sorting on propensity to coauthor: Implications for academic promotion. \emph{Economic Inquiry} \textbf{30}, 68--82 (1992).}

\bibitem[\citeproctext]{ref-holman2019researchers}
\CSLLeftMargin{16. }%DIF > 
\CSLRightInline{L. Holman, C. Morandin, Researchers collaborate with same-gendered colleagues more often than expected across the life sciences. \emph{PloS one} \textbf{14}, e0216128 (2019).}

\bibitem[\citeproctext]{ref-boschini2007team}
\CSLLeftMargin{17. }%DIF > 
\CSLRightInline{A. Boschini, A. Sjögren, Is team formation gender neutral? Evidence from coauthorship patterns. \emph{Journal of Labor Economics} \textbf{25}, 325--365 (2007).}

\bibitem[\citeproctext]{ref-kwiek2021gender}
\CSLLeftMargin{18. }%DIF > 
\CSLRightInline{M. Kwiek, W. Roszka, Gender-based homophily in research: A large-scale study of man-woman collaboration. \emph{Journal of Informetrics} \textbf{15}, 101171 (2021).}

\bibitem[\citeproctext]{ref-nguyen2025mind}
\CSLLeftMargin{19. }%DIF > 
\CSLRightInline{A. D. Nguyen, J. Jiang, B. Keller, Mind the gap: Gender, homophily and the glass ceiling in academic networks. \emph{Proceedings of the ACM on Human-Computer Interaction} \textbf{9}, 1--25 (2025).}

\bibitem[\citeproctext]{ref-ductor2023gender}
\CSLLeftMargin{20. }%DIF > 
\CSLRightInline{L. Ductor, S. Goyal, A. Prummer, Gender and collaboration. \emph{Review of Economics and Statistics} \textbf{105}, 1366--1378 (2023).}

\bibitem[\citeproctext]{ref-dion2018gendered}
\CSLLeftMargin{21. }%DIF > 
\CSLRightInline{M. L. Dion, J. L. Sumner, S. M. Mitchell, Gendered citation patterns across political science and social science methodology fields. \emph{Political analysis} \textbf{26}, 312--327 (2018).}

\bibitem[\citeproctext]{ref-mishra2018self}
\CSLLeftMargin{22. }%DIF > 
\CSLRightInline{S. Mishra, B. D. Fegley, J. Diesner, V. I. Torvik, Self-citation is the hallmark of productive authors, of any gender. \emph{PloS one} \textbf{13}, e0195773 (2018).}

\bibitem[\citeproctext]{ref-card2022gender}
\CSLLeftMargin{23. }%DIF > 
\CSLRightInline{D. Card, S. DellaVigna, P. Funk, N. Iriberri, Gender differences in peer recognition by economists. \emph{Econometrica} \textbf{90}, 1937--1971 (2022).}

\bibitem[\citeproctext]{ref-boivin2024sexism}
\CSLLeftMargin{24. }%DIF > 
\CSLRightInline{N. Boivin, S. Täuber, U. Beisiegel, U. Keller, J. G. Hering, Sexism in academia is bad for science and a waste of public funding. \emph{Nature Reviews Materials} \textbf{9}, 1--3 (2024).}

\bibitem[\citeproctext]{ref-newman2009first}
\CSLLeftMargin{25. }%DIF > 
\CSLRightInline{M. E. Newman, The first-mover advantage in scientific publication. \emph{Europhysics letters} \textbf{86}, 68001 (2009).}

\bibitem[\citeproctext]{ref-harris2020diversify}
\CSLLeftMargin{26. }%DIF > 
\CSLRightInline{J. K. Harris, M. A. Croston, E. T. Hutti, A. A. Eyler, Diversify the syllabi: Underrepresentation of female authors in college course readings. \emph{PLoS One} \textbf{15}, e0239012 (2020).}

\bibitem[\citeproctext]{ref-colgan2017gender}
\CSLLeftMargin{27. }%DIF > 
\CSLRightInline{J. Colgan, Gender bias in international relations graduate education? New evidence from syllabi. \emph{PS: Political Science \& Politics} \textbf{50}, 456--460 (2017).}

\bibitem[\citeproctext]{ref-skitka2021gender}
\CSLLeftMargin{28. }%DIF > 
\CSLRightInline{L. J. Skitka, Z. J. Melton, A. B. Mueller, K. Y. Wei, The gender gap: Who is (and is not) included on graduate-level syllabi in social/personality psychology. \emph{Personality and Social Psychology Bulletin} \textbf{47}, 863--872 (2021).}

\bibitem[\citeproctext]{ref-smith2020gender}
\CSLLeftMargin{29. }%DIF > 
\CSLRightInline{A. E. Smith, H. Hardt, P. Meister, H. J. Kim, Gender, race, age, and national origin predict whether faculty assign female-authored readings in graduate syllabi. \emph{PS: Political Science \& Politics} \textbf{53}, 100--106 (2020).}

\bibitem[\citeproctext]{ref-pinheiro2022women}
\CSLLeftMargin{30. }%DIF > 
\CSLRightInline{H. Pinheiro, M. Durning, D. Campbell, Do women undertake interdisciplinary research more than men, and do self-citations bias observed differences? \emph{Quantitative science studies} \textbf{3}, 363--392 (2022).}

\bibitem[\citeproctext]{ref-specht2022interdisciplinary}
\CSLLeftMargin{31. }%DIF > 
\CSLRightInline{A. Specht, K. Crowston, Interdisciplinary collaboration from diverse science teams can produce significant outcomes. \emph{PLoS One} \textbf{17}, e0278043 (2022).}

\bibitem[\citeproctext]{ref-yang2022gender}
\CSLLeftMargin{32. }%DIF > 
\CSLRightInline{Y. Yang, T. Y. Tian, T. K. Woodruff, B. F. Jones, B. Uzzi, Gender-diverse teams produce more novel and higher-impact scientific ideas. \emph{Proceedings of the National Academy of Sciences} \textbf{119}, e2200841119 (2022).}

\bibitem[\citeproctext]{ref-apesteguia2012impact}
\CSLLeftMargin{33. }%DIF > 
\CSLRightInline{J. Apesteguia, G. Azmat, N. Iriberri, The impact of gender composition on team performance and decision making: Evidence from the field. \emph{Management Science} \textbf{58}, 78--93 (2012).}

\bibitem[\citeproctext]{ref-opstrup2015right}
\CSLLeftMargin{34. }%DIF > 
\CSLRightInline{N. Opstrup, A. R. Villadsen, The right mix? Gender diversity in top management teams and financial performance. \emph{Public Administration Review} \textbf{75}, 291--301 (2015).}

\bibitem[\citeproctext]{ref-bear2011role}
\CSLLeftMargin{35. }%DIF > 
\CSLRightInline{J. B. Bear, A. W. Woolley, The role of gender in team collaboration and performance. \emph{Interdisciplinary science reviews} \textbf{36}, 146--153 (2011).}

\bibitem[\citeproctext]{ref-campbell2013gender}
\CSLLeftMargin{36. }%DIF > 
\CSLRightInline{L. G. Campbell, S. Mehtani, M. E. Dozier, J. Rinehart, Gender-heterogeneous working groups produce higher quality science. \emph{PloS one} \textbf{8}, e79147 (2013).}

\bibitem[\citeproctext]{ref-love2022impact}
\CSLLeftMargin{37. }%DIF > 
\CSLRightInline{H. B. Love, A. Stephens, B. K. Fosdick, E. Tofany, E. R. Fisher, The impact of gender diversity on scientific research teams: A need to broaden and accelerate future research. \emph{Humanities and Social Sciences Communications} \textbf{9}, 1--12 (2022).}

\bibitem[\citeproctext]{ref-viglione2020women}
\CSLLeftMargin{38. }%DIF > 
\CSLRightInline{G. Viglione, \emph{et al.}, Are women publishing less during the pandemic? Here's what the data say. \emph{Nature} \textbf{581}, 365--366 (2020).}

\bibitem[\citeproctext]{ref-gao2021potentially}
\CSLLeftMargin{39. }%DIF > 
\CSLRightInline{J. Gao, Y. Yin, K. R. Myers, K. R. Lakhani, D. Wang, Potentially long-lasting effects of the pandemic on scientists. \emph{Nature communications} \textbf{12}, 6188 (2021).}

\bibitem[\citeproctext]{ref-cui2022gender}
\CSLLeftMargin{40. }%DIF > 
\CSLRightInline{R. Cui, H. Ding, F. Zhu, Gender inequality in research productivity during the COVID-19 pandemic. \emph{Manufacturing \& Service Operations Management} \textbf{24}, 707--726 (2022).}

\bibitem[\citeproctext]{ref-garcia2024silver}
\CSLLeftMargin{41. }%DIF > 
\CSLRightInline{D. Garcı́a-Costa, F. Grimaldo, G. Bravo, B. Mehmani, F. Squazzoni, The silver lining of COVID-19 restrictions: Research output of academics under lockdown. \emph{Scientometrics} \textbf{129}, 1771--1786 (2024).}

\bibitem[\citeproctext]{ref-venturini2025effectsremoteworkingscientific}
\CSLLeftMargin{42. }%DIF > 
\CSLRightInline{S. Venturini, \emph{et al.}, \href{https://arxiv.org/abs/2511.18481}{The effects of remote working on scientific collaboration and impact}. (2025).}

\bibitem[\citeproctext]{ref-olechnicka2025impact}
\CSLLeftMargin{43. }%DIF > 
\CSLRightInline{A. Olechnicka, A. Ploszaj, E. Zegler-Poleska, The impact of the virtualization of scholarly conferences on the gender structure of conference contributors. \emph{Scientometrics} \textbf{130}, 423--445 (2025).}

\bibitem[\citeproctext]{ref-bray2022together}
\CSLLeftMargin{44. }%DIF > 
\CSLRightInline{H. J. Bray, \emph{et al.}, Together alone: Going online during covid-19 is changing scientific conferences. \emph{Challenges} \textbf{13}, 7 (2022).}

\bibitem[\citeproctext]{ref-oecd2022}
\CSLLeftMargin{45. }%DIF > 
\CSLRightInline{OECD, {``\href{https://doi.org/10.1787/bb6ee273-en}{What is the profile of academic staff and what is the student-academic staff ratio?}''} in \emph{Education at a Glance 2022: OECD Indicators}, (OECD Publishing, 2022).}

\bibitem[\citeproctext]{ref-lin2023sciscinet}
\CSLLeftMargin{46. }%DIF > 
\CSLRightInline{Z. Lin, Y. Yin, L. Liu, D. Wang, SciSciNet: A large-scale open data lake for the science of science research. \emph{Scientific Data} \textbf{10}, 315 (2023).}

\bibitem[\citeproctext]{ref-uzzi2013atypical}
\CSLLeftMargin{47. }%DIF > 
\CSLRightInline{B. Uzzi, S. Mukherjee, M. Stringer, B. Jones, Atypical combinations and scientific impact. \emph{Science} \textbf{342}, 468--472 (2013).}

\bibitem[\citeproctext]{ref-aksnes2011female}
\CSLLeftMargin{48. }%DIF > 
\CSLRightInline{D. W. Aksnes, K. Rorstad, F. Piro, G. Sivertsen, Are female researchers less cited? A large-scale study of norwegian scientists. \emph{Journal of the American Society for Information Science and Technology} \textbf{62}, 628--636 (2011).}

\bibitem[\citeproctext]{ref-chatterjee2021gender}
\CSLLeftMargin{49. }%DIF > 
\CSLRightInline{P. Chatterjee, R. M. Werner, Gender disparity in citations in high-impact journal articles. \emph{JAMA network open} \textbf{4}, e2114509--e2114509 (2021).}

\bibitem[\citeproctext]{ref-tekles2022same}
\CSLLeftMargin{50. }%DIF > 
\CSLRightInline{A. Tekles, K. Auspurg, L. Bornmann, Same-gender citations do not indicate a substantial gender homophily bias. \emph{Plos one} \textbf{17}, e0274810 (2022).}

\bibitem[\citeproctext]{ref-evans2016measuring}
\CSLLeftMargin{51. }%DIF > 
\CSLRightInline{E. D. Evans, Measuring interdisciplinarity using text. \emph{Socius} \textbf{2}, 2378023116654147 (2016).}

\bibitem[\citeproctext]{ref-agresti2000simple}
\CSLLeftMargin{52. }%DIF > 
\CSLRightInline{A. Agresti, B. Caffo, Simple and effective confidence intervals for proportions and differences of proportions result from adding two successes and two failures. \emph{The American Statistician} \textbf{54}, 280--288 (2000).}

\bibitem[\citeproctext]{ref-guimera2005team}
\CSLLeftMargin{53. }%DIF > 
\CSLRightInline{R. Guimera, B. Uzzi, J. Spiro, L. A. N. Amaral, Team assembly mechanisms determine collaboration network structure and team performance. \emph{Science} \textbf{308}, 697--702 (2005).}

\bibitem[\citeproctext]{ref-maddi2021gender}
\CSLLeftMargin{54. }%DIF > 
\CSLRightInline{A. Maddi, Y. Gingras, Gender diversity in research teams and citation impact in economics and management. \emph{Journal of Economic Surveys} \textbf{35}, 1381--1404 (2021).}

\bibitem[\citeproctext]{ref-santamaria2018comparison}
\CSLLeftMargin{55. }%DIF > 
\CSLRightInline{L. Santamarı́a, H. Mihaljević, Comparison and benchmark of name-to-gender inference services. \emph{PeerJ Computer Science} \textbf{4}, e156 (2018).}

\bibitem[\citeproctext]{ref-teele2017gender}
\CSLLeftMargin{56. }%DIF > 
\CSLRightInline{D. L. Teele, K. Thelen, Gender in the journals: Publication patterns in political science. \emph{PS: Political Science \& Politics} \textbf{50}, 433--447 (2017).}

\end{CSLReferences}

\DIFaddend \newpage
\appendix
\renewcommand{\thefigure}{SI-\arabic{figure}}
\renewcommand{\thesection}{SI}
\renewcommand{\thesubsection}{SI-\arabic{subsection}}
\DIFaddbegin 

\DIFaddend \setcounter{figure}{0}
\DIFaddbegin \setcounter{table}{0}
\DIFaddend \setcounter{section}{0}
\setcounter{subsection}{0}
\renewcommand{\thetable}{SI-\arabic{table}}

\section{Supporting Information}\label{supporting-information}

\subsection{\DIFdelbegin \DIFdel{Simulating Gender-Neutral Courses}\DIFdelend \DIFaddbegin \DIFadd{Data}\DIFaddend }\DIFdelbegin %DIFDELCMD < \label{si-simulations}
%DIFDELCMD < %%%
\DIFdelend \DIFaddbegin \label{data}
\DIFaddend 

\DIFdelbegin \DIFdel{We employed a simlar methodology to that developed by (\mbox{%DIFAUXCMD
\citeproc{ref-uzzi2013atypical}{36}}\hskip0pt%DIFAUXCMD
) for the analysis of academic citations.We counted the frequency of gender combinations (male-male, female-male, etc.) of each syllabus per field and academic year. We compared these combinations against those expected by chance, using a gender team composition network. In thisnetwork, for a given field, institution, and academic year, we switched all the instructors using a Monte Carlo algorithm. The switching algorithm preserves the total gender counts and the distribution of team size. This ensures that a course with \(n\) instructors in the original data will have the same number of instructors in the randomised network. Similarly, an institution with \(m\) male instructors and \(f\) female instructors teaching in a specific field will have the same number of male and female instructors.The only difference between the randomised and the original data will be the gender composition of the teams.}\DIFdelend \DIFaddbegin \subsubsection{\DIFadd{Gender disambiguation}}\label{gender-disambiguation}
\DIFaddend 

\DIFdelbegin \subsection{\DIFdel{Interdiscipline Similarity}}%DIFAUXCMD
\addtocounter{subsection}{-1}%DIFAUXCMD
\DIFdelend \DIFaddbegin \DIFadd{OpenSyllabus uses a Python implementation of Genderize.io to estimate the probability that an author's name corresponds to a female gender.}\footnote{\DIFadd{See }\url{https://github.com/opensyllabus/osp-gender}} \DIFadd{Names with a predicted probability above 0.5 are classified as female; all others are classified as male. Prior research indicates that Genderize.io achieves high overall accuracy (92\%) (\mbox{%DIFAUXCMD
\citeproc{ref-santamaria2018comparison}{55}}\hskip0pt%DIFAUXCMD
). In the context of academic authorship, a study hand-checking gender predictions in scholarly publications found that Genderize.io has an error rate of only about 2\% (\mbox{%DIFAUXCMD
\citeproc{ref-teele2017gender}{56}}\hskip0pt%DIFAUXCMD
). Despite this reliability, many instructors listed in syllabi provide only initials (e.g., ``J.'' instead of ``John'' or ``Jane''), which introduces ambiguity and increases noise in gender classification. To mitigate this, OpenSyllabus assigns an ``unknown'' category when an author's first name cannot be reliably determined.
}

\subsubsection{\DIFadd{Course Level}}\label{course-level}

\DIFadd{To determine the level of each course listed in the syllabi, we first extracted the course codes for each syllabus (e.g., Psychology 101, Econ101). Then, we standardised all course codes by converting them to uppercase and removing any non-alphanumeric characters. From each cleaned code, we extracted the first sequence of digits and used the leading digit of this number as an indicator of course level. Codes beginning with 1 or 2 were classified as }\emph{\DIFadd{basic}}\DIFadd{, those beginning with 3 or 4 as }\emph{\DIFadd{advanced}}\DIFadd{, and those beginning with 5 or higher as }\emph{\DIFadd{graduate}}\DIFadd{. Codes that did not contain a recognisable numeric component were assigned an }\emph{\DIFadd{unknown}} \DIFadd{level. This approach enabled consistent categorisation across heterogeneous code formats. Furthermore, on a subset of 56,000 U.S. syllabi, prior work (\mbox{%DIFAUXCMD
\citeproc{ref-biasi2022education}{10}}\hskip0pt%DIFAUXCMD
) showed that the course levels obtained using this procedure achieve an accuracy of 85\% when compared with the official course levels listed in the corresponding catalog. This provides confidence in the quality of the classification for the U.S. context. For EU countries, no comparable benchmark currently exists.
}

\subsubsection{\DIFadd{Interdiscipline Similarity}}\DIFaddend \label{si-interdisciplinarity}

We measured interdisciplinarity using text similarity between syllabi as in (\DIFdelbegin \DIFdel{\mbox{%DIFAUXCMD
\citeproc{ref-evans2016measuring}{26}}\hskip0pt%DIFAUXCMD
}\DIFdelend \DIFaddbegin \DIFadd{\mbox{%DIFAUXCMD
\citeproc{ref-evans2016measuring}{51}}\hskip0pt%DIFAUXCMD
}\DIFaddend ) and (\citeproc{ref-han2023interdisciplinary}{3}). We transformed text from course descriptions into ``bags of words,'' with term frequences (TF) normalised using the inverse document frequency (IDF). For each year, we calculated the TF-IDF scores for all syllabi and for 69 academic fields, using concatenated descriptions for the fields. We then computed the weighted average of the cosine similarity, \(\cos (i, f)\), between each course \(i\) and field \(f\), where the weight is based on the similarity between field \(f\) and the course's closest field \(f_{\max}\). Specifically, \(w(f_{\max}, f) = \cos(f_{\max}, f)\), where \(f_{\max}\) is the field that has the highest cosine similarity with course \(i\) (i.e., \(f_{\max} = \arg\max_f \cos(i, f)\)) in that academic year.

Thus, the interdisciplinarity score for course \(i\) is:
\[
  \text{interdisciplinarity score}_i = 1 - \frac{\sum_{f \in (1, 69)} \cos(f_{\max}, f) \cdot \cos(i, f)}{\sum_{f \in (1, 69)} \cos(f_{\max}, f)}. 
\]
To reduce computational costs, we ``bootstrap'' the field-by-field cosine similarity matrix by using a 10\% random sample of syllabi for each academic year. We repeated the sub-sampling procedure ten times and averaged the results.

\DIFdelbegin %DIFDELCMD < \clearpage
%DIFDELCMD < %%%
\DIFdelend \DIFaddbegin \subsection{\DIFadd{Simulating Gender-Neutral Courses}}\label{si-simulations}
\DIFaddend 

\DIFaddbegin \DIFadd{We employed a simlar methodology to that developed by (\mbox{%DIFAUXCMD
\citeproc{ref-uzzi2013atypical}{47}}\hskip0pt%DIFAUXCMD
) for the analysis of academic citations. We counted the frequency of gender combinations (male-male, female-male, etc.) of each syllabus per field and academic year. We compared these combinations against those expected by chance, using a gender team composition network. In this network, for a given field, institution, and academic year, we switched all the instructors using a Monte Carlo algorithm. The switching algorithm preserves the total gender counts and the distribution of team size. This ensures that a course with \(n\) instructors in the original data will have the same number of instructors in the randomised network. Similarly, an institution with \(m\) male instructors and \(f\) female instructors teaching in a specific field will have the same number of male and female instructors. The only difference between the randomised and the original data will be the gender composition of the teams.
}

\DIFaddend \subsection{Addtional Tables}\label{addtional-tables}

\DIFaddbegin \DIFadd{In this section, we present additional descriptive statistics for the number of syllabi in each country \ref{tab:SI-country}, the distribution of syllabi per field \ref{tab:SI-fields}, academic year \ref{tab:SI-years}, and teaching teams \ref{tab:SI-composition}.
}

\DIFaddend 
\begin{table}
[!h]
\centering
\caption{\label{tab:SI-country}Number of Syllabi per Country}
\centering
\fontsize{10}{12}\selectfont
\DIFdelbeginFL %DIFDELCMD < \begin{tabular}[t]{lrr}
%DIFDELCMD < \toprule
%DIFDELCMD < Country & N (thousands) & \%\\
%DIFDELCMD < \midrule
%DIFDELCMD < \cellcolor{gray!10}{Austria} & \cellcolor{gray!10}{33.3} & \cellcolor{gray!10}{0.6}\\
%DIFDELCMD < Other & 55.3 & 1.0\\
%DIFDELCMD < \cellcolor{gray!10}{Canada} & \cellcolor{gray!10}{389.8} & \cellcolor{gray!10}{7.3}\\
%DIFDELCMD < Germany & 62.4 & 1.2\\
%DIFDELCMD < \cellcolor{gray!10}{Denmark} & \cellcolor{gray!10}{24.3} & \cellcolor{gray!10}{0.5}\\
%DIFDELCMD < Spain & 37.8 & 0.7\\
%DIFDELCMD < \cellcolor{gray!10}{France} & \cellcolor{gray!10}{5.9} & \cellcolor{gray!10}{0.1}\\
%DIFDELCMD < Great Britain & 583.6 & 10.9\\
%DIFDELCMD < \cellcolor{gray!10}{Ireland} & \cellcolor{gray!10}{50.9} & \cellcolor{gray!10}{0.9}\\
%DIFDELCMD < Italy & 299.9 & 5.6\\
%DIFDELCMD < \cellcolor{gray!10}{Netherlands} & \cellcolor{gray!10}{86.5} & \cellcolor{gray!10}{1.6}\\
%DIFDELCMD < Poland & 107.2 & 2.0\\
%DIFDELCMD < \cellcolor{gray!10}{Portugal} & \cellcolor{gray!10}{42.1} & \cellcolor{gray!10}{0.8}\\
%DIFDELCMD < Sweden & 42.1 & 0.8\\
%DIFDELCMD < \cellcolor{gray!10}{USA} & \cellcolor{gray!10}{3554.1} & \cellcolor{gray!10}{66.1}\\
%DIFDELCMD < \bottomrule
%DIFDELCMD < \end{tabular}
%DIFDELCMD < %%%
\DIFdelendFL \DIFaddbeginFL \begin{tabu} to \linewidth {>{\raggedright}X>{\raggedleft}X>{\raggedleft}X}
\toprule
Country & Syllabi (thousands) & \%\\
\midrule
\cellcolor{gray!10}{USA} & \cellcolor{gray!10}{3554.1} & \cellcolor{gray!10}{66.1}\\
Great Britain & 583.6 & 10.9\\
\cellcolor{gray!10}{Canada} & \cellcolor{gray!10}{389.8} & \cellcolor{gray!10}{7.3}\\
Italy & 299.9 & 5.6\\
\cellcolor{gray!10}{Poland} & \cellcolor{gray!10}{107.2} & \cellcolor{gray!10}{2.0}\\
Netherlands & 86.5 & 1.6\\
\cellcolor{gray!10}{Germany} & \cellcolor{gray!10}{62.4} & \cellcolor{gray!10}{1.2}\\
Other & 55.3 & 1.0\\
\cellcolor{gray!10}{Ireland} & \cellcolor{gray!10}{50.9} & \cellcolor{gray!10}{0.9}\\
Portugal & 42.1 & 0.8\\
\cellcolor{gray!10}{Sweden} & \cellcolor{gray!10}{42.1} & \cellcolor{gray!10}{0.8}\\
Spain & 37.8 & 0.7\\
\cellcolor{gray!10}{Austria} & \cellcolor{gray!10}{33.3} & \cellcolor{gray!10}{0.6}\\
Denmark & 24.3 & 0.5\\
\cellcolor{gray!10}{France} & \cellcolor{gray!10}{5.9} & \cellcolor{gray!10}{0.1}\\
\bottomrule
\end{tabu}
\DIFaddendFL \end{table}


\begin{table}
[!h]
\centering
\caption{\label{tab:SI-fields}Number of Syllabi per Field}
\centering
\fontsize{10}{12}\selectfont
\DIFdelbeginFL %DIFDELCMD < \begin{tabular}[t]{lrrlrr}
%DIFDELCMD < \toprule
%DIFDELCMD < Field & N (thousands) & \% & Field & N (thousands) & \%\\
%DIFDELCMD < \midrule
%DIFDELCMD < \cellcolor{gray!10}{Accounting} & \cellcolor{gray!10}{76.9} & \cellcolor{gray!10}{1.4} & \cellcolor{gray!10}{Agriculture} & \cellcolor{gray!10}{36.2} & \cellcolor{gray!10}{0.7}\\
%DIFDELCMD < Anthropology & 51.9 & 1 & Architecture & 48.5 & 0.9\\
%DIFDELCMD < \cellcolor{gray!10}{Astronomy} & \cellcolor{gray!10}{13.9} & \cellcolor{gray!10}{0.3} & \cellcolor{gray!10}{Atmospheric Sciences} & \cellcolor{gray!10}{6.6} & \cellcolor{gray!10}{0.1}\\
%DIFDELCMD < Basic Computer Skills & 51.1 & 0.9 & Basic Skills & 50.7 & 0.9\\
%DIFDELCMD < \cellcolor{gray!10}{Biology} & \cellcolor{gray!10}{250.3} & \cellcolor{gray!10}{4.7} & \cellcolor{gray!10}{Business} & \cellcolor{gray!10}{450.4} & \cellcolor{gray!10}{8.4}\\
%DIFDELCMD < Career Skills & 13.9 & 0.3 & Chemistry & 106.4 & 2\\
%DIFDELCMD < \cellcolor{gray!10}{Chinese} & \cellcolor{gray!10}{10.6} & \cellcolor{gray!10}{0.2} & \cellcolor{gray!10}{Classics} & \cellcolor{gray!10}{23.3} & \cellcolor{gray!10}{0.4}\\
%DIFDELCMD < Computer Science & 310.5 & 5.8 & Construction & 10.9 & 0.2\\
%DIFDELCMD < \cellcolor{gray!10}{Cosmetology} & \cellcolor{gray!10}{6.9} & \cellcolor{gray!10}{0.1} & \cellcolor{gray!10}{Criminal Justice} & \cellcolor{gray!10}{45.5} & \cellcolor{gray!10}{0.8}\\
%DIFDELCMD < Criminology & 10.5 & 0.2 & Culinary Arts & 9.7 & 0.2\\
%DIFDELCMD < \cellcolor{gray!10}{Dance} & \cellcolor{gray!10}{14.2} & \cellcolor{gray!10}{0.3} & \cellcolor{gray!10}{Dentistry} & \cellcolor{gray!10}{13.9} & \cellcolor{gray!10}{0.3}\\
%DIFDELCMD < Earth Sciences & 49.5 & 0.9 & Economics & 130.9 & 2.4\\
%DIFDELCMD < \cellcolor{gray!10}{Education} & \cellcolor{gray!10}{240.3} & \cellcolor{gray!10}{4.5} & \cellcolor{gray!10}{Engineering} & \cellcolor{gray!10}{215.1} & \cellcolor{gray!10}{4}\\
%DIFDELCMD < Engineering Technician & 61.2 & 1.1 & English Literature & 357.3 & 6.6\\
%DIFDELCMD < \cellcolor{gray!10}{Film and Photography} & \cellcolor{gray!10}{49.8} & \cellcolor{gray!10}{0.9} & \cellcolor{gray!10}{Fine Arts} & \cellcolor{gray!10}{112.4} & \cellcolor{gray!10}{2.1}\\
%DIFDELCMD < Fitness and Leisure & 93.0 & 1.7 & French & 33.5 & 0.6\\
%DIFDELCMD < \cellcolor{gray!10}{Geography} & \cellcolor{gray!10}{43.3} & \cellcolor{gray!10}{0.8} & \cellcolor{gray!10}{German} & \cellcolor{gray!10}{42.2} & \cellcolor{gray!10}{0.8}\\
%DIFDELCMD < Health Technician & 31.9 & 0.6 & Hebrew & 3.1 & 0.1\\
%DIFDELCMD < \cellcolor{gray!10}{History} & \cellcolor{gray!10}{227.4} & \cellcolor{gray!10}{4.2} & \cellcolor{gray!10}{Japanese} & \cellcolor{gray!10}{10.7} & \cellcolor{gray!10}{0.2}\\
%DIFDELCMD < Journalism & 19.9 & 0.4 & Law & 99.5 & 1.9\\
%DIFDELCMD < \cellcolor{gray!10}{Liberal Arts} & \cellcolor{gray!10}{16.7} & \cellcolor{gray!10}{0.3} & \cellcolor{gray!10}{Library Science} & \cellcolor{gray!10}{25.0} & \cellcolor{gray!10}{0.5}\\
%DIFDELCMD < Linguistics & 32.3 & 0.6 & Marketing & 51.6 & 1\\
%DIFDELCMD < \cellcolor{gray!10}{Mathematics} & \cellcolor{gray!10}{406.2} & \cellcolor{gray!10}{7.6} & \cellcolor{gray!10}{Mechanic / Repair Tech} & \cellcolor{gray!10}{23.4} & \cellcolor{gray!10}{0.4}\\
%DIFDELCMD < Media / Communications & 118.8 & 2.2 & Medicine & 110.0 & 2\\
%DIFDELCMD < \cellcolor{gray!10}{Military Science} & \cellcolor{gray!10}{2.9} & \cellcolor{gray!10}{0.1} & \cellcolor{gray!10}{Music} & \cellcolor{gray!10}{101.2} & \cellcolor{gray!10}{1.9}\\
%DIFDELCMD < Natural Resource Management & 28.0 & 0.5 & Nursing & 98.5 & 1.8\\
%DIFDELCMD < \cellcolor{gray!10}{Nutrition} & \cellcolor{gray!10}{19.6} & \cellcolor{gray!10}{0.4} & \cellcolor{gray!10}{Philosophy} & \cellcolor{gray!10}{70.8} & \cellcolor{gray!10}{1.3}\\
%DIFDELCMD < Physics & 80.4 & 1.5 & Political Science & 167.3 & 3.1\\
%DIFDELCMD < \cellcolor{gray!10}{Psychology} & \cellcolor{gray!10}{222.7} & \cellcolor{gray!10}{4.1} & \cellcolor{gray!10}{Public Administration} & \cellcolor{gray!10}{7.8} & \cellcolor{gray!10}{0.1}\\
%DIFDELCMD < Public Safety & 7.4 & 0.1 & Religion & 29.9 & 0.6\\
%DIFDELCMD < \cellcolor{gray!10}{Sign Language} & \cellcolor{gray!10}{6.9} & \cellcolor{gray!10}{0.1} & \cellcolor{gray!10}{Social Work} & \cellcolor{gray!10}{48.8} & \cellcolor{gray!10}{0.9}\\
%DIFDELCMD < Sociology & 107.5 & 2 & Spanish & 60.5 & 1.1\\
%DIFDELCMD < \cellcolor{gray!10}{Theatre Arts} & \cellcolor{gray!10}{36.3} & \cellcolor{gray!10}{0.7} & \cellcolor{gray!10}{Theology} & \cellcolor{gray!10}{36.9} & \cellcolor{gray!10}{0.7}\\
%DIFDELCMD < Transportation & 4.3 & 0.1 & Veterinary Medicine & 11.9 & 0.2\\
%DIFDELCMD < \cellcolor{gray!10}{Women's Studies} & \cellcolor{gray!10}{7.4} & \cellcolor{gray!10}{0.1} & \cellcolor{gray!10}{} & \cellcolor{gray!10}{} & \cellcolor{gray!10}{}\\
%DIFDELCMD < \bottomrule
%DIFDELCMD < \end{tabular}
%DIFDELCMD < %%%
\DIFdelendFL \DIFaddbeginFL \begin{tabu} to \linewidth {>{\raggedright}X>{\raggedleft}X>{\raggedleft}X>{\raggedright}X>{\raggedleft}X>{\raggedleft}X}
\toprule
Field & N (thousands) & \% & Field & N (thousands) & \%\\
\midrule
\cellcolor{gray!10}{Accounting} & \cellcolor{gray!10}{76.9} & \cellcolor{gray!10}{1.4} & \cellcolor{gray!10}{Agriculture} & \cellcolor{gray!10}{36.2} & \cellcolor{gray!10}{0.7}\\
Anthropology & 51.9 & 1 & Architecture & 48.5 & 0.9\\
\cellcolor{gray!10}{Astronomy} & \cellcolor{gray!10}{13.9} & \cellcolor{gray!10}{0.3} & \cellcolor{gray!10}{Atmospheric Sciences} & \cellcolor{gray!10}{6.6} & \cellcolor{gray!10}{0.1}\\
Basic Computer Skills & 51.1 & 0.9 & Basic Skills & 50.7 & 0.9\\
\cellcolor{gray!10}{Biology} & \cellcolor{gray!10}{250.3} & \cellcolor{gray!10}{4.7} & \cellcolor{gray!10}{Business} & \cellcolor{gray!10}{450.4} & \cellcolor{gray!10}{8.4}\\
Career Skills & 13.9 & 0.3 & Chemistry & 106.4 & 2\\
\cellcolor{gray!10}{Chinese} & \cellcolor{gray!10}{10.6} & \cellcolor{gray!10}{0.2} & \cellcolor{gray!10}{Classics} & \cellcolor{gray!10}{23.3} & \cellcolor{gray!10}{0.4}\\
Computer Science & 310.5 & 5.8 & Construction & 10.9 & 0.2\\
\cellcolor{gray!10}{Cosmetology} & \cellcolor{gray!10}{6.9} & \cellcolor{gray!10}{0.1} & \cellcolor{gray!10}{Criminal Justice} & \cellcolor{gray!10}{45.5} & \cellcolor{gray!10}{0.8}\\
Criminology & 10.5 & 0.2 & Culinary Arts & 9.7 & 0.2\\
\cellcolor{gray!10}{Dance} & \cellcolor{gray!10}{14.2} & \cellcolor{gray!10}{0.3} & \cellcolor{gray!10}{Dentistry} & \cellcolor{gray!10}{13.9} & \cellcolor{gray!10}{0.3}\\
Earth Sciences & 49.5 & 0.9 & Economics & 130.9 & 2.4\\
\cellcolor{gray!10}{Education} & \cellcolor{gray!10}{240.3} & \cellcolor{gray!10}{4.5} & \cellcolor{gray!10}{Engineering} & \cellcolor{gray!10}{215.1} & \cellcolor{gray!10}{4}\\
Engineering Technician & 61.2 & 1.1 & English Literature & 357.3 & 6.6\\
\cellcolor{gray!10}{Film and Photography} & \cellcolor{gray!10}{49.8} & \cellcolor{gray!10}{0.9} & \cellcolor{gray!10}{Fine Arts} & \cellcolor{gray!10}{112.4} & \cellcolor{gray!10}{2.1}\\
Fitness and Leisure & 93.0 & 1.7 & French & 33.5 & 0.6\\
\cellcolor{gray!10}{Geography} & \cellcolor{gray!10}{43.3} & \cellcolor{gray!10}{0.8} & \cellcolor{gray!10}{German} & \cellcolor{gray!10}{42.2} & \cellcolor{gray!10}{0.8}\\
Health Technician & 31.9 & 0.6 & Hebrew & 3.1 & 0.1\\
\cellcolor{gray!10}{History} & \cellcolor{gray!10}{227.4} & \cellcolor{gray!10}{4.2} & \cellcolor{gray!10}{Japanese} & \cellcolor{gray!10}{10.7} & \cellcolor{gray!10}{0.2}\\
Journalism & 19.9 & 0.4 & Law & 99.5 & 1.9\\
\cellcolor{gray!10}{Liberal Arts} & \cellcolor{gray!10}{16.7} & \cellcolor{gray!10}{0.3} & \cellcolor{gray!10}{Library Science} & \cellcolor{gray!10}{25.0} & \cellcolor{gray!10}{0.5}\\
Linguistics & 32.3 & 0.6 & Marketing & 51.6 & 1\\
\cellcolor{gray!10}{Mathematics} & \cellcolor{gray!10}{406.2} & \cellcolor{gray!10}{7.6} & \cellcolor{gray!10}{Mechanic / Repair Tech} & \cellcolor{gray!10}{23.4} & \cellcolor{gray!10}{0.4}\\
Media / Communications & 118.8 & 2.2 & Medicine & 110.0 & 2\\
\cellcolor{gray!10}{Military Science} & \cellcolor{gray!10}{2.9} & \cellcolor{gray!10}{0.1} & \cellcolor{gray!10}{Music} & \cellcolor{gray!10}{101.2} & \cellcolor{gray!10}{1.9}\\
Natural Resource Management & 28.0 & 0.5 & Nursing & 98.5 & 1.8\\
\cellcolor{gray!10}{Nutrition} & \cellcolor{gray!10}{19.6} & \cellcolor{gray!10}{0.4} & \cellcolor{gray!10}{Philosophy} & \cellcolor{gray!10}{70.8} & \cellcolor{gray!10}{1.3}\\
Physics & 80.4 & 1.5 & Political Science & 167.3 & 3.1\\
\cellcolor{gray!10}{Psychology} & \cellcolor{gray!10}{222.7} & \cellcolor{gray!10}{4.1} & \cellcolor{gray!10}{Public Administration} & \cellcolor{gray!10}{7.8} & \cellcolor{gray!10}{0.1}\\
Public Safety & 7.4 & 0.1 & Religion & 29.9 & 0.6\\
\cellcolor{gray!10}{Sign Language} & \cellcolor{gray!10}{6.9} & \cellcolor{gray!10}{0.1} & \cellcolor{gray!10}{Social Work} & \cellcolor{gray!10}{48.8} & \cellcolor{gray!10}{0.9}\\
Sociology & 107.5 & 2 & Spanish & 60.5 & 1.1\\
\cellcolor{gray!10}{Theatre Arts} & \cellcolor{gray!10}{36.3} & \cellcolor{gray!10}{0.7} & \cellcolor{gray!10}{Theology} & \cellcolor{gray!10}{36.9} & \cellcolor{gray!10}{0.7}\\
Transportation & 4.3 & 0.1 & Veterinary Medicine & 11.9 & 0.2\\
\cellcolor{gray!10}{Women's Studies} & \cellcolor{gray!10}{7.4} & \cellcolor{gray!10}{0.1} & \cellcolor{gray!10}{} & \cellcolor{gray!10}{} & \cellcolor{gray!10}{}\\
\bottomrule
\end{tabu}
\DIFaddendFL \end{table}


\begin{table}
[!h]
\centering
\caption{\label{tab:SI-years}Number of Syllabi per Year}
\centering
\fontsize{10}{12}\selectfont
\DIFdelbeginFL %DIFDELCMD < \begin{tabular}[t]{lrr}
%DIFDELCMD < \toprule
%DIFDELCMD < Academic year & N (thousands) & \%\\
%DIFDELCMD < \midrule
%DIFDELCMD < \cellcolor{gray!10}{1999 or older} & \cellcolor{gray!10}{40.5} & \cellcolor{gray!10}{0.8}\\
%DIFDELCMD < 2000 & 20.0 & 0.4\\
%DIFDELCMD < \cellcolor{gray!10}{2001} & \cellcolor{gray!10}{27.2} & \cellcolor{gray!10}{0.5}\\
%DIFDELCMD < 2002 & 40.1 & 0.7\\
%DIFDELCMD < \cellcolor{gray!10}{2003} & \cellcolor{gray!10}{55.1} & \cellcolor{gray!10}{1.0}\\
%DIFDELCMD < 2004 & 73.8 & 1.4\\
%DIFDELCMD < \cellcolor{gray!10}{2005} & \cellcolor{gray!10}{77.2} & \cellcolor{gray!10}{1.4}\\
%DIFDELCMD < 2006 & 95.5 & 1.8\\
%DIFDELCMD < \cellcolor{gray!10}{2007} & \cellcolor{gray!10}{100.6} & \cellcolor{gray!10}{1.9}\\
%DIFDELCMD < 2008 & 124.0 & 2.3\\
%DIFDELCMD < \cellcolor{gray!10}{2009} & \cellcolor{gray!10}{153.0} & \cellcolor{gray!10}{2.8}\\
%DIFDELCMD < 2010 & 200.4 & 3.7\\
%DIFDELCMD < \cellcolor{gray!10}{2011} & \cellcolor{gray!10}{256.3} & \cellcolor{gray!10}{4.8}\\
%DIFDELCMD < 2012 & 302.2 & 5.6\\
%DIFDELCMD < \cellcolor{gray!10}{2013} & \cellcolor{gray!10}{375.3} & \cellcolor{gray!10}{7.0}\\
%DIFDELCMD < 2014 & 446.5 & 8.3\\
%DIFDELCMD < \cellcolor{gray!10}{2015} & \cellcolor{gray!10}{464.6} & \cellcolor{gray!10}{8.6}\\
%DIFDELCMD < 2016 & 547.8 & 10.2\\
%DIFDELCMD < \cellcolor{gray!10}{2017} & \cellcolor{gray!10}{658.5} & \cellcolor{gray!10}{12.3}\\
%DIFDELCMD < 2018 & 669.2 & 12.5\\
%DIFDELCMD < \cellcolor{gray!10}{2019} & \cellcolor{gray!10}{647.0} & \cellcolor{gray!10}{12.0}\\
%DIFDELCMD < \bottomrule
%DIFDELCMD < \end{tabular}
%DIFDELCMD < %%%
\DIFdelendFL \DIFaddbeginFL \begin{tabu} to \linewidth {>{\raggedright}X>{\raggedleft}X>{\raggedleft}X}
\toprule
Academic year & N (thousands) & \%\\
\midrule
\cellcolor{gray!10}{1999 or older} & \cellcolor{gray!10}{40.5} & \cellcolor{gray!10}{0.8}\\
2000 & 20.0 & 0.4\\
\cellcolor{gray!10}{2001} & \cellcolor{gray!10}{27.2} & \cellcolor{gray!10}{0.5}\\
2002 & 40.1 & 0.7\\
\cellcolor{gray!10}{2003} & \cellcolor{gray!10}{55.1} & \cellcolor{gray!10}{1.0}\\
2004 & 73.8 & 1.4\\
\cellcolor{gray!10}{2005} & \cellcolor{gray!10}{77.2} & \cellcolor{gray!10}{1.4}\\
2006 & 95.5 & 1.8\\
\cellcolor{gray!10}{2007} & \cellcolor{gray!10}{100.6} & \cellcolor{gray!10}{1.9}\\
2008 & 124.0 & 2.3\\
\cellcolor{gray!10}{2009} & \cellcolor{gray!10}{153.0} & \cellcolor{gray!10}{2.8}\\
2010 & 200.4 & 3.7\\
\cellcolor{gray!10}{2011} & \cellcolor{gray!10}{256.3} & \cellcolor{gray!10}{4.8}\\
2012 & 302.2 & 5.6\\
\cellcolor{gray!10}{2013} & \cellcolor{gray!10}{375.3} & \cellcolor{gray!10}{7.0}\\
2014 & 446.5 & 8.3\\
\cellcolor{gray!10}{2015} & \cellcolor{gray!10}{464.6} & \cellcolor{gray!10}{8.6}\\
2016 & 547.8 & 10.2\\
\cellcolor{gray!10}{2017} & \cellcolor{gray!10}{658.5} & \cellcolor{gray!10}{12.3}\\
2018 & 669.2 & 12.5\\
\cellcolor{gray!10}{2019} & \cellcolor{gray!10}{647.0} & \cellcolor{gray!10}{12.0}\\
\bottomrule
\end{tabu}
\DIFaddendFL \end{table}


\begin{table}
[!h]
\centering
\caption{\label{tab:SI-composition}Teaching Team Configurations}
\centering
\fontsize{10}{12}\selectfont
\DIFdelbeginFL %DIFDELCMD < \begin{tabular}[t]{lrr}
%DIFDELCMD < \toprule
%DIFDELCMD < Team composition & N (thousands) & \%\\
%DIFDELCMD < \midrule
%DIFDELCMD < \cellcolor{gray!10}{Male only (M)} & \cellcolor{gray!10}{2587.6} & \cellcolor{gray!10}{48.1}\\
%DIFDELCMD < Female only (F) & 1828.6 & 34.0\\
%DIFDELCMD < \cellcolor{gray!10}{Male-Male (MM)} & \cellcolor{gray!10}{349.6} & \cellcolor{gray!10}{6.5}\\
%DIFDELCMD < Female-Female (FF) & 228.6 & 4.3\\
%DIFDELCMD < \cellcolor{gray!10}{Mixed Gender (MF/FM)} & \cellcolor{gray!10}{83.4} & \cellcolor{gray!10}{1.6}\\
%DIFDELCMD < Mixed Gender (MF/FM) & 67.9 & 1.3\\
%DIFDELCMD < \bottomrule
%DIFDELCMD < \end{tabular}
%DIFDELCMD < %%%
\DIFdelendFL \DIFaddbeginFL \begin{tabu} to \linewidth {>{\raggedright}X>{\raggedleft}X>{\raggedleft}X}
\toprule
Team composition & N (thousands) & \%\\
\midrule
\cellcolor{gray!10}{Male only (M)} & \cellcolor{gray!10}{2587.6} & \cellcolor{gray!10}{48.1}\\
Female only (F) & 1828.6 & 34.0\\
\cellcolor{gray!10}{Male-Male (MM)} & \cellcolor{gray!10}{349.6} & \cellcolor{gray!10}{6.5}\\
Female-Female (FF) & 228.6 & 4.3\\
\cellcolor{gray!10}{Mixed Gender (MF/FM)} & \cellcolor{gray!10}{83.4} & \cellcolor{gray!10}{1.6}\\
Mixed Gender (MF/FM) & 67.9 & 1.3\\
\bottomrule
\end{tabu}
\DIFaddendFL 
\end{table}

\DIFdelbegin %DIFDELCMD < \processdelayedfloats
%DIFDELCMD < \clearpage
%DIFDELCMD < \singlespacing
%DIFDELCMD < %%%
\DIFdelend \DIFaddbegin \subsection{\DIFadd{Representativeness of the gender distribution}}\label{representativeness-of-the-gender-distribution}
\DIFaddend 

\DIFdelbegin \section*{\DIFdel{References}}%DIFAUXCMD
%DIFDELCMD < \label{references}
%DIFDELCMD < %%%
\addcontentsline{toc}{section}{\DIFdel{References}}
%DIFAUXCMD
\DIFdelend \DIFaddbegin \DIFadd{One potential concern for our analysis is the representativeness of the data. Our dataset covers several million courses and thousands of educational institutions, corresponding to about 50\% of the educational organisations in the ROR, which makes the results interesting in their own right. Nevertheless, one may question how representative these data are of the broader population of instructors. Although we cannot directly test representativeness, we evaluate it in two ways. First, we examine whether the overall female-to-male ratio over time in each country correlates with OECD statistics on academic staff by gender. Second, for the U.S., we assess whether the sex ratio by field and year correlates with NCSES data on Ph.D.~graduates by academic field, country, and year, which serves as a proxy for the maximum number of employable instructors.
}\DIFaddend 

\DIFdelbegin %DIFDELCMD < \phantomsection\label{refs}
%DIFDELCMD < \begin{CSLReferences}{0}{1}
%DIFDELCMD < \bibitem[\citeproctext]{ref-bakkenes2010teacher}
%DIFDELCMD < \CSLLeftMargin{1. }%%%
%DIF < 
%DIFDELCMD < \CSLRightInline{I. Bakkenes, J. D. Vermunt, T. Wubbels, Teacher learning in the context of educational innovation: Learning activities and learning outcomes of experienced teachers. \emph{Learning and instruction} \textbf{20}, 533--548 (2010).}
%DIFDELCMD < %%%
\DIFdelend \DIFaddbegin \DIFadd{As illustrated in Figure \ref{fig:country-coverage}, our analysis of overall trends reveals a very high correlation for English-speaking countries (U.S. = .98, U.K. = .90), and greater heterogeneity for non-English-speaking countries, where representativeness is expected to be weaker given our focus on English-language syllabi. The analysis based on Ph.D.~graduates similarly shows strong field-level correlations in many fields (Figure \ref{fig:field-coverage-us}), such as in Natural Sciences (rho = .91), Social sciences and journalism (rho = .90), Education (rho = .77). However, there are fields with considerably lower correlation, such as in Health and welfare (rho = .14) or IT (rho = .29). This heterogeneity may stem from differences across disciplines in the proportion of Ph.D.~graduates who remain in academia (e.g., biology has a comparatively high share of graduates who pursue careers outside academia). Nevertheless, these patterns suggest that, at least for English-speaking countries and in many academic fields, our data are reasonably representative.
}\DIFaddend 

\DIFdelbegin \bibitem[\citeproctext]{ref-vembye2024effects}
%DIFDELCMD < \CSLLeftMargin{2. }%%%
%DIF < 
%DIFDELCMD < \CSLRightInline{M. H. Vembye, F. Weiss, B. Hamilton Bhat, The effects of co-teaching and related collaborative models of instruction on student achievement: A systematic review and meta-analysis. \emph{Review of Educational Research} \textbf{94}, 376--422 (2024).}
%DIFDELCMD < %%%
\DIFdelend \DIFaddbegin \begin{figure}
\centering
\pandocbounded{\includegraphics[keepaspectratio]{docs/report_files/figure-latex/country-coverage-1.pdf}}
\caption{\label{fig:country-coverage}\DIFaddFL{Women-men ratio of instructors in OpenSyllabus (OS) vs academic staff in tertiary education (ISCED 5-8) per country and year. Pearson correlation between OS and Eurostat series shown in each panel. Source: Eurostat (educ\_uoe\_perd02).}}

\end{figure}
\DIFaddend 

\DIFdelbegin \bibitem[\citeproctext]{ref-han2023interdisciplinary}
%DIFDELCMD < \CSLLeftMargin{3. }%%%
%DIF < 
%DIFDELCMD < \CSLRightInline{S. Han, J. LaViolette, C. Borkenhagen, W. McAllister, P. S. Bearman, Interdisciplinary college curriculum and its labor market implications. \emph{Proceedings of the National Academy of Sciences} \textbf{120}, e2221915120 (2023).}
%DIFDELCMD < %%%
\DIFdelend \DIFaddbegin \begin{figure}
\centering
\pandocbounded{\includegraphics[keepaspectratio]{docs/report_files/figure-latex/field-coverage-us-1.pdf}}
\caption{\label{fig:field-coverage-us}\DIFaddFL{Women-men ratio of instructors in OpenSyllabus (OS) vs PhD graduates in the US by academic field per year. Pearson correlation between OS and representative series shown in each panel. Source: NCES (Survey of Earned Doctorates).}}

\end{figure}
\DIFaddend 

\DIFdelbegin \bibitem[\citeproctext]{ref-ellemers2018gender}
%DIFDELCMD < \CSLLeftMargin{4. }%%%
%DIF < 
%DIFDELCMD < \CSLRightInline{N. Ellemers, Gender stereotypes. \emph{Annual review of psychology} \textbf{69}, 275--298 (2018).}
%DIFDELCMD < %%%
\DIFdelend \DIFaddbegin \subsection{\DIFadd{Pre-/Post- Covid-19 Patterns}}\label{pre-post--covid-19-patterns}
\DIFaddend 

\DIFdelbegin \bibitem[\citeproctext]{ref-lockwood2006someone}
%DIFDELCMD < \CSLLeftMargin{5. }%%%
%DIF < 
%DIFDELCMD < \CSLRightInline{P. Lockwood, {``Someone like me can be successful''}: Do college students need same-gender role models? \emph{Psychology of women quarterly} \textbf{30}, 36--46 (2006).}
%DIFDELCMD < %%%
\DIFdelend \DIFaddbegin \DIFadd{Figure \ref{fig:covid-trend} illustrates pre-/post-pandemic patterns in the teams across academic fields.
}\DIFaddend 

\DIFdelbegin \bibitem[\citeproctext]{ref-canaan2023impact}
%DIFDELCMD < \CSLLeftMargin{6. }%%%
%DIF < 
%DIFDELCMD < \CSLRightInline{S. Canaan, P. Mouganie, The impact of advisor gender on female students' STEM enrollment and persistence. \emph{Journal of Human Resources} \textbf{58}, 593--632 (2023).}
%DIFDELCMD < %%%
\DIFdelend \DIFaddbegin \begin{figure}
\centering
\pandocbounded{\includegraphics[keepaspectratio]{docs/report_files/figure-latex/covid-trend-1.pdf}}
\caption{\label{fig:covid-trend}\textbf{\DIFaddFL{Pre-/Post-Covid Courses}}\DIFaddFL{. Number of courses per CIP field (gray) in thirty-three European countries before and after Covid-19. Data from Open Syllabus (2017-2022). Smooth trend curves show the average pattern per broad field definition (ISCED-F). Most fields exhibit a significant drop in courses during the pandemic, especially Health and Welfare, Arts and Humanities, and Social Sciences and Journalism.}}

\end{figure}
\DIFaddend 

\DIFdelbegin \bibitem[\citeproctext]{ref-buckingham2021learning}
%DIFDELCMD < \CSLLeftMargin{7. }%%%
%DIF < 
%DIFDELCMD < \CSLRightInline{L. R. Buckingham, A. López-Hernández, B. Strotmann, Learning by comparison: The benefits of co-teaching for university professors' professional development in \emph{Frontiers in Education}, (Frontiers Media SA, 2021), p. 776991.}
%DIFDELCMD < %%%
\DIFdelend \DIFaddbegin \subsection{\DIFadd{Heterogeneity: US vs non-US}}\label{heterogeneity-us-vs-non-us}
\DIFaddend 

\DIFdelbegin \bibitem[\citeproctext]{ref-nowogrodzki2016mining}
%DIFDELCMD < \CSLLeftMargin{8. }%%%
%DIF < 
%DIFDELCMD < \CSLRightInline{A. Nowogrodzki, Mining the secrets of college syllabuses. \emph{Nature} \textbf{539}, 125--126 (2016).}
%DIFDELCMD < %%%
\DIFdelend \DIFaddbegin \DIFadd{Figure \ref{fig:hetero-us-plots} illustrates heterogeneity in the associations between teams and outcomes within US courses and the other countries.
}\DIFaddend 

\DIFdelbegin \bibitem[\citeproctext]{ref-gao2023quantifying}
%DIFDELCMD < \CSLLeftMargin{9. }%%%
%DIF < 
%DIFDELCMD < \CSLRightInline{J. Gao, D. Wang, Quantifying the benefit of artificial intelligence for scientific research. \emph{arXiv preprint arXiv:2304.10578} (2023).}
%DIFDELCMD < %%%
\DIFdelend \DIFaddbegin \begin{figure}
\centering
\pandocbounded{\includegraphics[keepaspectratio]{docs/report_files/figure-latex/hetero-us-plots-1.pdf}}
\caption{\label{fig:hetero-us-plots}\textbf{\DIFaddFL{Heterogeneity within US courses}}\DIFaddFL{. Estimates and 95\% confidence intervals for differences in (1) interdisciplinarity (2) readings' conventionality (3) age of readings (4) female-authored ratio and (5) readings' atypicality. Estimates are from separate mixed-effects regressions for courses from the US vs the full sample, in each academic year, controlling for country, field, institution, course level, probability that the document is a syllabus, and number of readings, with weights proportional to the square root of the number of readings. Smoothed curves show local estimates averaged across years, weighted by the precision (standard errors) of the estimates. The comparison reveals no noticeable differences.}}

\end{figure}
\DIFaddend 

\DIFdelbegin \bibitem[\citeproctext]{ref-biasi2022education}
%DIFDELCMD < \CSLLeftMargin{10. }%%%
%DIF < 
%DIFDELCMD < \CSLRightInline{B. Biasi, S. Ma, {``The education-innovation gap''} (National Bureau of Economic Research, 2022).}
%DIFDELCMD < %%%
\DIFdelend \DIFaddbegin \subsection{\DIFadd{Robustness Analysis}}\label{robustness-analysis}
\DIFaddend 

\DIFdelbegin \bibitem[\citeproctext]{ref-guimera2005team}
%DIFDELCMD < \CSLLeftMargin{11. }%%%
%DIF < 
%DIFDELCMD < \CSLRightInline{R. Guimera, B. Uzzi, J. Spiro, L. A. N. Amaral, Team assembly mechanisms determine collaboration network structure and team performance. \emph{Science} \textbf{308}, 697--702 (2005).}
%DIFDELCMD < %%%
\DIFdelend \DIFaddbegin \DIFadd{To ensure the reliability and robustness of our results, we conducted several additional analyses. In addition to reporting results from separate regressions on subsets of data by year, we also estimated mixed-effects regressions using the full dataset. These results are broadly consistent with those presented in the paper. Figure \ref{fig:robust-interdisc} shows robustness for interdisciplinarity, Figure \ref{fig:robust-women} for the fraction of female authors cited in the readings, Figures \ref{fig:robust-age-readings}, \ref{fig:robust-conventional}, \ref{fig:robust-atypical} show full-model results for the age, conventionality, and the atypicality of readings.
}\DIFaddend 

\DIFdelbegin \bibitem[\citeproctext]{ref-wuchty2007increasing}
%DIFDELCMD < \CSLLeftMargin{12. }%%%
%DIF < 
%DIFDELCMD < \CSLRightInline{S. Wuchty, B. F. Jones, B. Uzzi, The increasing dominance of teams in production of knowledge. \emph{Science} \textbf{316}, 1036--1039 (2007).}
%DIFDELCMD < %%%
\DIFdelend \DIFaddbegin \DIFadd{While we report differences in percentile rank per year to address skewness in the outcome variables, we also explored alternative specificationss. First, we log-transformed our dependent variables, including the age, uncononventionality and atypicalty of readings, to explore potential non-linear relationships, particularly relevant for disciplines with old readings (e.g., history) or based on citations. Second, we employed a quasi-Poisson mixed-effect model instead of the reported linear mixed-effect regression to better account for non-linearities in the count of female authors (Figure \ref{fig:robust-women}, panel B). These changes yield similar results.
}\DIFaddend 

\DIFdelbegin \bibitem[\citeproctext]{ref-milojevic2014principles}
%DIFDELCMD < \CSLLeftMargin{13. }%%%
%DIF < 
%DIFDELCMD < \CSLRightInline{S. Milojević, Principles of scientific research team formation and evolution. \emph{Proceedings of the National Academy of Sciences} \textbf{111}, 3984--3989 (2014).}
%DIFDELCMD < %%%
\DIFdelend \DIFaddbegin \DIFadd{Finally, we further explored differences in team formation that distinguish between male-led and female-led courses, using the gender of the first listed instructor as a proxy for seniority. Overall, team formation trends appear robust to these additional analyses (Figure \ref{fig:order-instructors}).
}\DIFaddend 

\DIFdelbegin \bibitem[\citeproctext]{ref-yang2022gender}
%DIFDELCMD < \CSLLeftMargin{14. }%%%
%DIF < 
%DIFDELCMD < \CSLRightInline{Y. Yang, T. Y. Tian, T. K. Woodruff, B. F. Jones, B. Uzzi, Gender-diverse teams produce more novel and higher-impact scientific ideas. \emph{Proceedings of the National Academy of Sciences} \textbf{119}, e2200841119 (2022).}
%DIFDELCMD < %%%
\DIFdelend \DIFaddbegin \begin{figure}
\centering
\pandocbounded{\includegraphics[keepaspectratio]{docs/report_files/figure-latex/robust-interdisc-1.pdf}}
\caption{\label{fig:robust-interdisc}\textbf{\DIFaddFL{Robustness analysis for interdisciplinarity.}} \DIFaddFL{Estimates and 95\% confidence intervals from a multiple linear regression on the full sample. Dependent variable = percentile rank per year of interdisciplinarity score. Controls include the country, year, the probability of the document being a proper syllabus, field of study, institution, course levels, and (logarithm of) the number of matched readings.}}

\end{figure}
\DIFaddend 

\DIFdelbegin \bibitem[\citeproctext]{ref-boschini2007team}
%DIFDELCMD < \CSLLeftMargin{15. }%%%
%DIF < 
%DIFDELCMD < \CSLRightInline{A. Boschini, A. Sjögren, Is team formation gender neutral? Evidence from coauthorship patterns. \emph{Journal of Labor Economics} \textbf{25}, 325--365 (2007).}
%DIFDELCMD < %%%
\DIFdelend \DIFaddbegin \begin{figure}
\centering
\pandocbounded{\includegraphics[keepaspectratio]{docs/report_files/figure-latex/robust-women-1.pdf}}
\caption{\label{fig:robust-women}\textbf{\DIFaddFL{Robustness analysis for female-authored readings.}} \DIFaddFL{Estimates and 95\% confidence intervals from (A) a multiple linear regression (B) quasi-poisson regression on the full sample. Dependent variable = fraction of female-authored readings to toal readings. Controls include the country, year, the probability of the document being a proper syllabus, field of study, institution, course levels, and (logarithm of) the number of matched readings.}}

\end{figure}
\DIFaddend 

\DIFdelbegin \bibitem[\citeproctext]{ref-love2022impact}
%DIFDELCMD < \CSLLeftMargin{16. }%%%
%DIF < 
%DIFDELCMD < \CSLRightInline{H. B. Love, A. Stephens, B. K. Fosdick, E. Tofany, E. R. Fisher, The impact of gender diversity on scientific research teams: A need to broaden and accelerate future research. \emph{Humanities and Social Sciences Communications} \textbf{9}, 1--12 (2022).}
%DIFDELCMD < %%%
\DIFdelend \DIFaddbegin \begin{figure}
\centering
\pandocbounded{\includegraphics[keepaspectratio]{docs/report_files/figure-latex/robust-age-readings-1.pdf}}
\caption{\label{fig:robust-age-readings}\textbf{\DIFaddFL{Robustness analysis for age of readings.}} \DIFaddFL{Estimates and 95\% confidence intervals from a multiple linear regression on the full sample. Dependent variable = percentile rank of the age of publication of a syllabus' readings as (A) percentile rank per year (B) logarithm. Controls include the country, year, the probability of the document being a proper syllabus, field of study, institution, course levels, and (logarithm of) the number of matched readings.}}

\end{figure}
\DIFaddend 

\DIFdelbegin \bibitem[\citeproctext]{ref-kwiek2021gender}
%DIFDELCMD < \CSLLeftMargin{17. }%%%
%DIF < 
%DIFDELCMD < \CSLRightInline{M. Kwiek, W. Roszka, Gender-based homophily in research: A large-scale study of man-woman collaboration. \emph{Journal of Informetrics} \textbf{15}, 101171 (2021).}
%DIFDELCMD < %%%
\DIFdelend \DIFaddbegin \begin{figure}
\centering
\pandocbounded{\includegraphics[keepaspectratio]{docs/report_files/figure-latex/robust-conventional-1.pdf}}
\caption{\label{fig:robust-conventional}\textbf{\DIFaddFL{Robustness analysis for conventionality of readings.}} \DIFaddFL{Estimates and 95\% confidence intervals from a multiple linear regression on the full sample. Dependent variable = median citations of a syllabus' readings as from references in previous syllabi as (A) percentile rank per year (B) logarithm. Controls include the country, year, the probability of the document being a proper syllabus, field of study, institution, course levels, and (logarithm of) the number of matched readings.}}

\end{figure}
\DIFaddend 

\DIFdelbegin \bibitem[\citeproctext]{ref-boivin2024sexism}
%DIFDELCMD < \CSLLeftMargin{18. }%%%
%DIF < 
%DIFDELCMD < \CSLRightInline{N. Boivin, S. Täuber, U. Beisiegel, U. Keller, J. G. Hering, Sexism in academia is bad for science and a waste of public funding. \emph{Nature Reviews Materials} \textbf{9}, 1--3 (2024).}
%DIFDELCMD < %%%
\DIFdelend \DIFaddbegin \begin{figure}
\centering
\pandocbounded{\includegraphics[keepaspectratio]{docs/report_files/figure-latex/robust-atypical-1.pdf}}
\caption{\label{fig:robust-atypical}\textbf{\DIFaddFL{Robustness analysis for atypicality of readings.}} \DIFaddFL{Estimates and 95\% confidence intervals from a multiple linear regression on the full sample. Dependent variable = median Uzzi et al's atypicality of a syllabus' readings as precomputed by SciSciNet v2 as (A) percentile rank per year (B) logarithm. Controls include the country, year, the probability of the document being a proper syllabus, field of study, institution, course levels, and (logarithm of) the number of matched readings.}}

\end{figure}
\DIFaddend 

\DIFdelbegin \bibitem[\citeproctext]{ref-weisshaar2017publish}
%DIFDELCMD < \CSLLeftMargin{19. }%%%
%DIF < 
%DIFDELCMD < \CSLRightInline{K. Weisshaar, Publish and perish? An assessment of gender gaps in promotion to tenure in academia. \emph{Social Forces} \textbf{96}, 529--560 (2017).}
%DIFDELCMD < 

%DIFDELCMD < \bibitem[\citeproctext]{ref-bedi2012gender}
%DIFDELCMD < \CSLLeftMargin{20. }%%%
%DIF < 
%DIFDELCMD < \CSLRightInline{G. Bedi, N. T. Van Dam, M. Munafo, Gender inequality in awarded research grants. \emph{The Lancet} \textbf{380}, 474 (2012).}
%DIFDELCMD < 

%DIFDELCMD < \bibitem[\citeproctext]{ref-sarsons2017recognition}
%DIFDELCMD < \CSLLeftMargin{21. }%%%
%DIF < 
%DIFDELCMD < \CSLRightInline{H. Sarsons, Recognition for group work: Gender differences in academia. \emph{American Economic Review} \textbf{107}, 141--145 (2017).}
%DIFDELCMD < 

%DIFDELCMD < \bibitem[\citeproctext]{ref-card2022gender}
%DIFDELCMD < \CSLLeftMargin{22. }%%%
%DIF < 
%DIFDELCMD < \CSLRightInline{D. Card, S. DellaVigna, P. Funk, N. Iriberri, Gender differences in peer recognition by economists. \emph{Econometrica} \textbf{90}, 1937--1971 (2022).}
%DIFDELCMD < 

%DIFDELCMD < \bibitem[\citeproctext]{ref-card2022impact}
%DIFDELCMD < \CSLLeftMargin{23. }%%%
%DIF < 
%DIFDELCMD < \CSLRightInline{D. Card, C. Domnisoru, S. G. Sanders, L. Taylor, V. Udalova, {``The impact of female teachers on female students' lifetime well-being''} (National Bureau of Economic Research, 2022).}
%DIFDELCMD < 

%DIFDELCMD < \bibitem[\citeproctext]{ref-bromham2016interdisciplinary}
%DIFDELCMD < \CSLLeftMargin{24. }%%%
%DIF < 
%DIFDELCMD < \CSLRightInline{L. Bromham, R. Dinnage, X. Hua, Interdisciplinary research has consistently lower funding success. \emph{Nature} \textbf{534}, 684--687 (2016).}
%DIFDELCMD < 

%DIFDELCMD < \bibitem[\citeproctext]{ref-lariviere2010relationship}
%DIFDELCMD < \CSLLeftMargin{25. }%%%
%DIF < 
%DIFDELCMD < \CSLRightInline{V. Larivière, Y. Gingras, On the relationship between interdisciplinarity and scientific impact. \emph{Journal of the American Society for Information Science and Technology} \textbf{61}, 126--131 (2010).}
%DIFDELCMD < 

%DIFDELCMD < \bibitem[\citeproctext]{ref-evans2016measuring}
%DIFDELCMD < \CSLLeftMargin{26. }%%%
%DIF < 
%DIFDELCMD < \CSLRightInline{E. D. Evans, Measuring interdisciplinarity using text. \emph{Socius} \textbf{2}, 2378023116654147 (2016).}
%DIFDELCMD < 

%DIFDELCMD < \bibitem[\citeproctext]{ref-skitka2021gender}
%DIFDELCMD < \CSLLeftMargin{27. }%%%
%DIF < 
%DIFDELCMD < \CSLRightInline{L. J. Skitka, Z. J. Melton, A. B. Mueller, K. Y. Wei, The gender gap: Who is (and is not) included on graduate-level syllabi in social/personality psychology. \emph{Personality and Social Psychology Bulletin} \textbf{47}, 863--872 (2021).}
%DIFDELCMD < 

%DIFDELCMD < \bibitem[\citeproctext]{ref-colgan2017gender}
%DIFDELCMD < \CSLLeftMargin{28. }%%%
%DIF < 
%DIFDELCMD < \CSLRightInline{J. Colgan, Gender bias in international relations graduate education? New evidence from syllabi. \emph{PS: Political Science \& Politics} \textbf{50}, 456--460 (2017).}
%DIFDELCMD < 

%DIFDELCMD < \bibitem[\citeproctext]{ref-harris2020diversify}
%DIFDELCMD < \CSLLeftMargin{29. }%%%
%DIF < 
%DIFDELCMD < \CSLRightInline{J. K. Harris, M. A. Croston, E. T. Hutti, A. A. Eyler, Diversify the syllabi: Underrepresentation of female authors in college course readings. \emph{PLoS One} \textbf{15}, e0239012 (2020).}
%DIFDELCMD < 

%DIFDELCMD < \bibitem[\citeproctext]{ref-bear2011role}
%DIFDELCMD < \CSLLeftMargin{30. }%%%
%DIF < 
%DIFDELCMD < \CSLRightInline{J. B. Bear, A. W. Woolley, The role of gender in team collaboration and performance. \emph{Interdisciplinary science reviews} \textbf{36}, 146--153 (2011).}
%DIFDELCMD < 

%DIFDELCMD < \bibitem[\citeproctext]{ref-maddi2021gender}
%DIFDELCMD < \CSLLeftMargin{31. }%%%
%DIF < 
%DIFDELCMD < \CSLRightInline{A. Maddi, Y. Gingras, Gender diversity in research teams and citation impact in economics and management. \emph{Journal of Economic Surveys} \textbf{35}, 1381--1404 (2021).}
%DIFDELCMD < 

%DIFDELCMD < \bibitem[\citeproctext]{ref-oecd2022}
%DIFDELCMD < \CSLLeftMargin{32. }%%%
%DIF < 
%DIFDELCMD < \CSLRightInline{OECD, {``\href{https://doi.org/10.1787/bb6ee273-en}{What is the profile of academic staff and what is the student-academic staff ratio?}''} in \emph{Education at a Glance 2022: OECD Indicators}, (OECD Publishing, 2022).}
%DIFDELCMD < 

%DIFDELCMD < \bibitem[\citeproctext]{ref-aksnes2011female}
%DIFDELCMD < \CSLLeftMargin{33. }%%%
%DIF < 
%DIFDELCMD < \CSLRightInline{D. W. Aksnes, K. Rorstad, F. Piro, G. Sivertsen, Are female researchers less cited? A large-scale study of norwegian scientists. \emph{Journal of the American Society for Information Science and Technology} \textbf{62}, 628--636 (2011).}
%DIFDELCMD < 

%DIFDELCMD < \bibitem[\citeproctext]{ref-chatterjee2021gender}
%DIFDELCMD < \CSLLeftMargin{34. }%%%
%DIF < 
%DIFDELCMD < \CSLRightInline{P. Chatterjee, R. M. Werner, Gender disparity in citations in high-impact journal articles. \emph{JAMA network open} \textbf{4}, e2114509--e2114509 (2021).}
%DIFDELCMD < 

%DIFDELCMD < \bibitem[\citeproctext]{ref-tekles2022same}
%DIFDELCMD < \CSLLeftMargin{35. }%%%
%DIF < 
%DIFDELCMD < \CSLRightInline{A. Tekles, K. Auspurg, L. Bornmann, Same-gender citations do not indicate a substantial gender homophily bias. \emph{Plos one} \textbf{17}, e0274810 (2022).}
%DIFDELCMD < 

%DIFDELCMD < \bibitem[\citeproctext]{ref-uzzi2013atypical}
%DIFDELCMD < \CSLLeftMargin{36. }%%%
%DIF < 
%DIFDELCMD < \CSLRightInline{B. Uzzi, S. Mukherjee, M. Stringer, B. Jones, Atypical combinations and scientific impact. \emph{Science} \textbf{342}, 468--472 (2013).}
%DIFDELCMD < 

%DIFDELCMD < \bibitem[\citeproctext]{ref-wang2017bias}
%DIFDELCMD < \CSLLeftMargin{37. }%%%
%DIF < 
%DIFDELCMD < \CSLRightInline{J. Wang, R. Veugelers, P. Stephan, Bias against novelty in science: A cautionary tale for users of bibliometric indicators. \emph{Research Policy} \textbf{46}, 1416--1436 (2017).}
%DIFDELCMD < 

%DIFDELCMD < \bibitem[\citeproctext]{ref-agresti2000simple}
%DIFDELCMD < \CSLLeftMargin{38. }%%%
%DIF < 
%DIFDELCMD < \CSLRightInline{A. Agresti, B. Caffo, Simple and effective confidence intervals for proportions and differences of proportions result from adding two successes and two failures. \emph{The American Statistician} \textbf{54}, 280--288 (2000).}
%DIFDELCMD < 

%DIFDELCMD < \bibitem[\citeproctext]{ref-pinheiro2022women}
%DIFDELCMD < \CSLLeftMargin{39. }%%%
%DIF < 
%DIFDELCMD < \CSLRightInline{H. Pinheiro, M. Durning, D. Campbell, Do women undertake interdisciplinary research more than men, and do self-citations bias observed differences? \emph{Quantitative science studies} \textbf{3}, 363--392 (2022).}
%DIFDELCMD < 

%DIFDELCMD < \end{CSLReferences}
%DIFDELCMD < %%%
\DIFdelend \DIFaddbegin \begin{figure}
\centering
\pandocbounded{\includegraphics[keepaspectratio]{docs/report_files/figure-latex/order-instructors-1.pdf}}
\caption{\label{fig:order-instructors}\textbf{\DIFaddFL{Robustness in the order of listed instructors}}\DIFaddFL{. Comparing the gap between simulated and actual mixed-gender teams when the first listed instructor is a woman versus when the first listed instructor is a man.}}

\end{figure}
\DIFaddend 

\end{document}
